%%%%%%%%%%%%%%%%%%%%%%%%%%%%%%%%%%%%%%%%%%%%%%%%%%%%%%%%%%%%%%%%%%%%%%
% Lớp 11 | Dãy số. Cấp số cộng. Cấp số nhân | Dãy số
% Dạng 1D2?1-2: Số hạng tổng quát, biểu diễn dãy số
% Số câu: 129  (sinh tự động bởi tach_cau.py; từ nay thầy sửa trực tiếp file này)
%%%%%%%%%%%%%%%%%%%%%%%%%%%%%%%%%%%%%%%%%%%%%%%%%%%%%%%%%%%%%%%%%%%%%%

%%%%%%%%%%%%%%%% TRẮC NGHIỆM 4 PHƯƠNG ÁN %%%%%%%%%%%%%%%%

% nguồn: Ok_Toan11_TracNghiem_Lan2
\begin{ex}%[1D2N1-2]%[ma:deffa23b]
	Cho dãy số  $\left(u_n \right)$, biết công thức số hạng tổng quát $u_n=2n-3$, $n\in\mathbb{N}^*$. Số hạng thứ $10$ của dãy số bằng
	\choice 
		 {\True $17$}
		 {$20$}
		{$10$} 
		 {$7$}
	\loigiai{ 
		Ta  có $ u_{10}=2\cdot10-3=17$.
	} 
\end{ex}

% nguồn: Ok_Toan11_TracNghiem_Lan2
\begin{ex}%[1D2N1-2]%[ma:3c2ea399]
	Cho dãy số $\left(u_n\right)$, biết $u_n=\dfrac{n}{2^n-1 }$, $n\in\mathbb{N}^{*}$. Ba số  hạng đầu tiên của dãy số là
	 \choice
		{$\dfrac{1}{2}$;  $\dfrac{2}{3}$; $\dfrac{1}{16}$}
		 {$1$; $\dfrac{1}{2}$; $\dfrac{3}{4}$}
		{$1$; $\dfrac{1}{4}$;  $\dfrac{1}{8}$}
		{\True $1$;  $\dfrac{2}{3}$; $\dfrac{3}{7}$}
	\loigiai{ 
		Ta có $u_1=1$, $u_2=\dfrac{2}{3}$,  $u_3=\dfrac{3}{7}$.
	 }
\end{ex}

% nguồn: Ok_Toan11_TracNghiem_Lan2
\begin{ex}%[1D2N1-2]%[ma:e73a640a]
	Cho  dãy số $\left(u_n \right)$ với $u_n=2^n$, $n\in \mathbb{N}^{*}$. Tìm số hạng $u_{n+1}$.
	 \choice
		{\True  $u_{n+1}=2^n\cdot2$}
		 {$u_{n+1}=2^n+1$}
		{$u_{n+1}=2\left(n+1\right)$} 
		 {$u_{n+1}=2^n+2$}
	 \loigiai{
		Ta  có $u_{n+1}=2^{n+1}=2\cdot2^n$.
	} 
\end{ex}

% nguồn: Ok_Toan11_TracNghiem_Lan2
\begin{ex}%[1D2N1-2]%[ma:f7fe6fc5]
	 ]
	Cho dãy số $\left(u_n \right)$, biết $u_n=\dfrac{2n+5}{5n-4}$, $n\in \mathbb{N}^{*}$. Số $\dfrac{7}{12}$ là  số hạng thứ mấy của dãy số?
	\choice 
		 {$6$}
		{\True $8$} 
		 {$9$}
		 {$10$}
	\loigiai{ 
		Ta có  $u_n=\dfrac{7}{12} \Leftrightarrow \dfrac{2n+5}{5n-4}=\dfrac{7}{12} \left(n\in \mathbb{N}^{*}\right)\Leftrightarrow 24n+60=35n-28\Leftrightarrow 11n=88\Leftrightarrow n=8$.
	 }
\end{ex}

% nguồn: Ok_Toan11_TracNghiem_Lan2
\begin{ex}%[1D2N1-2]%[ma:dd208c7b]
	Cho dãy số  $\left(u_n \right)$, biết $u_n=\dfrac{n+1}{2n+1}$, $n\in \mathbb{N}^{*}$. Số $\dfrac{8}{15}$ là số hạng thứ mấy của dãy số?
	\choice 
		 {$8$}
		{$6$} 
		{$5$} 
		{\True  $7$}
	\loigiai{ 
		Ta có $u_n=\dfrac{8}{15}  \Leftrightarrow \dfrac{n+1}{2n+1}=\dfrac{8}{15}\left(n\in \mathbb{N}^{*}\right) \Leftrightarrow 15n+15=16n+8\Leftrightarrow n=7$.
	 }
\end{ex}

% nguồn: Ok_Toan11_TracNghiem_Lan2
\begin{ex}%[1D2N1-2]%[ma:4b49d290]
	Cho dãy số $\left(u_n  \right)$ xác định bởi $\heva{&{u_1=\dfrac{1}{2}} \\
	&{u_n=\dfrac{1}{2-u_{n-1}},\ 2\le  n\in\mathbb{Z}}
	}$. Khi đó  $u_3$ có giá trị bằng
	\choice 
		{\True $\dfrac{3}{4}$} 
		{$\dfrac{4}{3}$} 
		{$\dfrac{2}{3}$} 
		{$\dfrac{3}{2}$} 
	 \loigiai{
		Theo công thức truy hồi ta  có $u_2=\dfrac{1}{2-\dfrac{1}{2}}=\dfrac{2}{3}$ $\Rightarrow u_3=\dfrac{1}{2-\dfrac{2}{3}}=\dfrac{3}{4}$.
	} 
\end{ex}

% nguồn: Ok_Toan11_TracNghiem_Lan2
\begin{ex}%[1D2N1-2]%[ma:317659d5]
	Cho dãy số $\left(u_n \right)$ với $u_n=5^{n+1}$, $n\in \mathbb{N}^{*}$. Tìm  số hạng $u_{n-1}$.
	 \choice
		{$u_{n-1}=5^{n-1}$} 
		{\True $u_{n-1}=5^n$} 
		 {$u_{n-1}=5\cdot5^{n+1}$}
		 {$u_{n-1}=5\cdot5^{n-1}$}
	\loigiai{ 
		Ta có  $u_{n-1}=5^{\left(n-1\right)+1}=5^n$.
	 }
\end{ex}

% nguồn: Ok_Toan11_TracNghiem_Lan2
\begin{ex}%[1D2N1-2]%[ma:f173957c]
	Cho dãy số $\left(u_n \right)$ xác định bởi $\heva{&{u_1=-2} \\ 
	 &{u_n=2u_{n-1}+n^2,\ 2\le n\in\mathbb{Z}}}$. Số hạng thứ tư của dãy số đó~bằng
	 \choice
		 {$0$}
		{$93$} 
		 {$9$}
		 {\True $34$}
	\loigiai{ 
		Ta có $u_2=2\cdot\left(-2\right)+2^2=0$, $u_3=2\cdot0+3^3=9$,  $u_4=2\cdot9+4^2=34$.
	} 
\end{ex}

% nguồn: Ok_Toan11_TracNghiem_Lan2
\begin{ex}%[1D2N1-2]%[ma:f1b060d0]
	Cho dãy  số $\left(u_n \right)$ có $u_n=\dfrac{n^2-1}{n^2+1}$, $n\in\mathbb{N}^*$. Tính $u_2$.
	 \choice
		 {$u_2=\dfrac{1}{5}$}
		 {$u_2=\dfrac{2}{5}$}
		{\True  $u_2=\dfrac{3}{5}$}
		{$u_2=\dfrac{4}{5}$} 
	 \loigiai{
		 Ta có $u_2=\dfrac{2^2-1}{2^2+1}=\dfrac{3}{5}$.
	 }
\end{ex}

% nguồn: Ok_Toan11_TracNghiem_Lan2
\begin{ex}%[1D2N1-2]%[ma:c7d2ca23]
	Cho dãy số $\left(u_n \right)$ có công thức số hạng tổng quát $u_n=8-3n$, $n\in\mathbb{N}^*$. Tính $u_4$. 
	 \choice
		 {$2$}
		 {$-7$}
		 {$-5$}
		{\True $-4$} 
	 \loigiai{
		Ta có  $u_4=8-3\cdot4=-4$.
	 }
\end{ex}

% nguồn: Ok_Toan11_TracNghiem_Lan2
\begin{ex}%[1D2N1-2]%[ma:98fa5ae3]
	Cho dãy số $\left(u_n \right)$, biết $u_n=2\cdot3^n$, $n\in\mathbb{N}^*$. Giá trị  của $u_{20}$ bằng
	\choice 
		{$2\cdot3^{19}$} 
		{\True  $2\cdot3^{20}$}
		{$3^{20}$} 
		{$2\cdot3^{21\;}$} 
	\loigiai{ 
		Ta  có $u_{20}=2\cdot3^{20}$.
	} 
\end{ex}

% nguồn: Ok_Toan11_TracNghiem_Lan2
\begin{ex}%[1D2N1-2]%[ma:75c40b66]
	Cho dãy số $\left(u_n \right)$ với $u_n=3^n$, $n\in  \mathbb{N}^{*}$. Số hạng $u_{n+1}$ bằng
	 \choice
		 {$3^n+1$}
		 {$3^n+3$}
		 {\True $3^n\cdot3$}
		 {$3(n+1)$}
	\loigiai{ 
		 Ta có $u_{n+1}=3^{n+1}=3^n\cdot3$.
	 }
\end{ex}

% nguồn: Ok_Toan11_TracNghiem_Lan2
\begin{ex}%[1D2N1-2]%[ma:f078dada]
	Cho dãy số $\left(u_n \right)$ với $u_n=\left(\dfrac{n-1}{n+1} \right)^{2n+3}$, $n\in  \mathbb{N}^{*}$. Tìm số hạng $u_{n+1}$.
	\choice 
		{$u_{n+1}=\left(\dfrac{n-1}{n+1} \right)^{2\left(n+1\right)+3}$} 
		{$u_{n+1}=\left(\dfrac{n-1}{n+1} \right)^{2\left(n-1\right)+3}$} 
		 {$u_{n+1}=\left(\dfrac{n}{n+2} \right)^{2n+3}$}
		{\True $u_{n+1}=\left(\dfrac{n}{n+2} \right)^{2n+5}$} 
	 \loigiai{
		Ta có $u_{n+1}=\left(\dfrac{n+1-1}{n+1+1}  \right)^{2\left(n+1\right)+3}=\left(\dfrac{n}{n+2} \right)^{2n+5}$.
	} 
\end{ex}

% nguồn: Ok_Toan11_TracNghiem_Lan2
\begin{ex}%[1D2N1-2]%[ma:fe1455c6]
	Cho  dãy số $\left(u_n \right)$ xác định bởi $u_n=\dfrac{n-1}{n^2+2n+3}$, $n\in\mathbb{N}^*$. Giá trị $u_{21}$ là
	\choice 
		 {$\dfrac{11}{243}$}
		{\True  $\dfrac{10}{243}$}
		 {$\dfrac{21}{443}$}
		{$\dfrac{19}{443}$} 
	\loigiai{ 
		 Ta có $u_{21}=\dfrac{21-1}{21^2+2\cdot21+3}=\dfrac{10}{243}$.
	} 
\end{ex}

% nguồn: Ok_Toan11_TracNghiem_Lan2
\begin{ex}%[1D2N1-2]%[ma:559160ef]
	Cho  dãy số $\left(u_n \right)$ xác định bởi $\heva{&{u_1=2} \\
	&{u_{n+1}=\dfrac{1}{3} \left(u_n+1\right),\ n\in\mathbb{N}^{*}.}}$ Tìm  số hạng $u_4$.
	\choice 
		{\True $u_4=\dfrac{5}{9}$} 
		{$u_4=1$} 
		 {$u_4=\dfrac{2}{3}$}
		 {$u_4=\dfrac{14}{27}$}
	\loigiai{ 
		Ta có $u_2=\dfrac{1}{3}  \left(2+1\right)=1$, $ u_3=\dfrac{1}{3} \left(1+1\right)=\dfrac{2}{3}$, $ u_4=\dfrac{1}{3} \left(\dfrac{2}{3}+1\right)=\dfrac{5}{9}$.
	} 
\end{ex}

% nguồn: Ok_Toan11_TracNghiem_Lan2
\begin{ex}%[1D2N1-2]%[ma:3ec24fcb]
	 Cho dãy số $\left(u_n \right)$ với $u_n=3^n$, $n\in \mathbb{N}^{*}$. Tìm số hạng $u_{2n-1}$.
	 \choice
		 {$u_{2n-1}=3^2\cdot3^n-1$}
		{\True $u_{2n-1}=3^n\cdot3^{n-1}$} 
		{$u_{2n-1}=3^{2n}-1$} 
		 {$u_{2n-1}=3^{2\left(n-1\right)}$}
	 \loigiai{
		Ta có $u_{2n-1}=3^{2n-1}=3^n\cdot3^{n-1}$. 
	} 
\end{ex}

% nguồn: Ok_Toan11_TracNghiem_Lan2
\begin{ex}%[1D2N1-2]%[ma:6b32e966]
	Cho dãy số $\left(u_n \right)$, biết $u_n=\dfrac{n-1}{n^2+1}$, $n\in \mathbb{N}^{*}$. Số $\dfrac{2}{13}$ là số  hạng thứ mấy của dãy số?
	\choice 
		 {Thứ ba}
		 {Thứ tư}
		{\True Thứ năm} 
		{Thứ sáu} 
	 \loigiai{
		Ta  có
		\begin{center}
			$u_n=\dfrac{2}{13} \Leftrightarrow \dfrac{n-1}{n^2+1}=\dfrac{2}{13}\left(n\in  \mathbb{N}^{*}\right) \Leftrightarrow 13n-13=2n^2+2\Leftrightarrow 2n^2-13n+15=0\Leftrightarrow \hoac{&{n=5\left(\text{nhận}\right)} \\
			 &{n=\dfrac{3}{2}\left(\text{loại}\right).}
			}$ 
		\end{center}
	} 
\end{ex}

% nguồn: Ok_Toan11_TracNghiem_Lan2
\begin{ex}%[1D2N1-2]%[ma:2067c0a2]
	Cho  dãy số $\left(u_n \right)$ xác định bởi $\heva{&{u_1=7} \\
	&{u_{n+1}=2u_n+3,\  n\in\mathbb{N}^{*}}
	}$. Khi đó  $u_5$ bằng
	 \choice
		 {$317$}
		{\True $157$} 
		{$77$} 
		 {$112$}
	\loigiai{ 
		Ta có $u_2=2\cdot7+3=17$,  $u_3=2\cdot17+3=37$, $u_4=2\cdot37+3=77$, $u_5=2\cdot77+3=157$.
	} 
\end{ex}

% nguồn: Ok_Toan11_TracNghiem_Lan2
\begin{ex}%[1D2N1-2]%[ma:ed04caba]
	Cho dãy số $\left(u_n \right)$, biết $u_n=n^3-8n^2-5n+7$, $n\in  \mathbb{N}^{*}$. Số $-33$ là số hạng thứ mấy của dãy số?
	 \choice
		{$5$} 
		 {$6$}
		{\True $8$} 
		{$9$} 
	 \loigiai{
		Ta  có
		\begin{center}
			$u_n=-33\Leftrightarrow  n^3-8n^2-5n+7=-33\Leftrightarrow n^3-8n^2-5n+40=0\Leftrightarrow \hoac{&{n=8\left(\text{nhận}\right)} \\
			&{n=\pm \sqrt{5}  \left(\text{loại}\right).}
			 }$
		\end{center}
	} 
\end{ex}

% nguồn: Ok_Toan11_TracNghiem_Lan2
\begin{ex}%[1D2N1-2]%[ma:1c38a094]
	Cho dãy số $\left(u_n  \right)$ với $u_n=2n+3$, $n\in\mathbb{N}^*$. Tìm số hạng thứ $6$ của dãy số.
	 \choice
		{$17$} 
		 {$5$}
		 {\True $15$}
		{$7$} 
	\loigiai{ 
		Ta có số hạng thứ  $6$ của dãy là $u_6=2\cdot6+3=15$.
	} 
\end{ex}

% nguồn: Ok_Toan11_TracNghiem_Lan2
\begin{ex}%[1D2N1-2]%[ma:6c3dc835]
	Cho dãy số $\left(u_n \right)$ xác định bởi $\heva{&{u_1=3} \\ 
	&{u_{n+1}=\dfrac{u_n}{2}+2,\ n\in\mathbb{N}^{*}} 
	}$. Mệnh  đề nào sau đây \textbf{sai}?
	 \choice
		 {\True $u_2=\dfrac{5}{2}$}
		{$u_3=\dfrac{15}{4}$} 
		 {$u_4=\dfrac{31}{8}$}
		 {$u_5=\dfrac{63}{16}$}
	 \loigiai{
		Vì $u_2=\dfrac{3}{2}+2=\dfrac{7}{2}$. 
	 }
\end{ex}

% nguồn: Ok_Toan11_TracNghiem_Lan2
\begin{ex}%[1D2N1-2]%[ma:b5536e0e]
	Cho dãy số $\left(u_n \right)$ xác định  bởi $\heva{&{u_1=-1} \\
	 &{u_{n+1}=u_n+3,\ n\in\mathbb{N}^{*}}}$. Ba số hạng đầu tiên của dãy số đó là
	\choice 
		 {\True $-1$; $2$; $5$}
		{$1$;  $4$; $7$}
		{$4$; $7$;  $10$}
		 {$-1$; $3$; $7$}
	\loigiai{ 
		 Ta có $u_1=-1$, $u_2=-1+3=2$, $u_3=2+3=5$.
	} 
\end{ex}

% nguồn: Ok_Toan11_TracNghiem_Lan2
\begin{ex}%[1D2N1-2]%[ma:9eef2e4e]
	 Cho dãy số $\left(u_n \right)$ xác định bởi $\heva{&{u_1=3} \\
	&{u_{n+1}=u_n+5,\ n\in\mathbb{N}^{*}}}$.  Ba số hạng đầu tiên của dãy số đó là
	\choice 
		{$-3$; $6$;  $9$}
		{$3$; $-2$;  $-7$}
		{\True  $3$; $8$; $13$}
		{$3$;  $5$; $7$}
	\loigiai{ 
		Ta  có $u_1=3$, $u_2=3+5=8$, $u_3=8+5=13$.
	 }
\end{ex}

% nguồn: Ok_Toan11_TracNghiem_Lan2
\begin{ex}%[1D2N1-2]%[ma:cc37f345]
	Cho dãy số $\left(u_n \right)$ với $u_n=3^n$, $n\in \mathbb{N}^{*}$. Số hạng  $u_{2n}$ bằng
	\choice 
		 {$3^n+3$}
		 {\True $9^n$}
		 {$3^n\cdot3$}
		{$4^{2n}$} 
	 \loigiai{
		 Ta có $u_{2n}=3^{2n}=9^n$.
	} 
\end{ex}

% nguồn: Ok_TracNghiemLop11_Lan1
\begin{ex}%[1D2N1-2]%[ma:27d5210e]
	Cho dãy số $(u_n)$ có công thức số hạng tổng  quát $u_n=3^n-10,n\in{\mathbb{N}^{*}}$, tính giá trị $u_3$.
	 \choice
		{\True $u_3=17$} 
		 {$u_3=27$}
		 {$u_3=-1$}
		 {$u_3=-10$}
	\loigiai{ 
	Ta  có $u_3=3^3-10=17$. }
\end{ex}

% nguồn: Ok_TracNghiemLop11_Lan1
\begin{ex}%[1D2N1-2]%[ma:d08b40bf]
	Cho dãy  số $\left(u_n\right)$ có các số hạng đầu là $\dfrac{1}{5}$, $\dfrac{1}{5^2}$, $\dfrac{1}{5^3}$, $\dfrac{1}{5^4}$, $\dfrac{1}{5^5}$, $\cdots $ Số hạng tổng quát của dãy số này là
	 \choice
		 {$u_n=\dfrac{1}{5}.\dfrac{1}{5^{n+1}}$}
		 {$u_n=\dfrac{1}{5^{n+1}}$}
		 {\True $u_n=\dfrac{1}{5^n}$}
		 {$u_n=\dfrac{1}{5^{n-1}}$}
	\loigiai{ 
		Dãy $u_n=\dfrac{1}{5^n}$ có khai khiển là $\dfrac{1}{5}$, $\dfrac{1}{5^2}$, $\dfrac{1}{5^3}$,  $\dfrac{1}{5^4}$, $\dfrac{1}{5^5}$, $\cdots$
	 }
\end{ex}

% nguồn: Ok_TracNghiemLop11_Lan1
\begin{ex}%[1D2N1-2]%[ma:51ec1dc1]
	Cho dãy số $\left(u_n\right)$ với  $u_n=11^n.$ Tìm số hạng $u_{n+1}$.
	 \choice
		{\True $u_{n+1}=11\cdot11^n$} 
		{$u_{n+1}=11^n+11$} 
		 {$u_{n+1}=11\left(n+1\right)$}
		 {$u_{n+1}=11^n+1$}
	\loigiai{ 
		 Ta có
	$u_{n+1}=11^{n+1}=11\cdot11^n.$} 
\end{ex}

% nguồn: Ok_Toan11_TracNghiem_Lan2
\begin{ex}%[1D2H1-2]%[ma:164177e9]
	Cho dãy số  $(u_n)$ với $\heva{&u_1=2 \\ &u_{n+1}=2u_n, \; n \ge 2}$. Công thức số hạng tổng quát của dãy số này là
	\choice 
		{$u_n=n^{n-1}$} 
		{\True  $u_n=2^n$}
		{$u_n=2^{n+1}$} 
		{$u_n=2$} 
	 \loigiai{
		 Ta có
		$\heva{&u_1=2 \\ 
		&u_2=2u_1 \\ 
		&u_3=2u_2  \\
		&\vdots \\ 
		&u_n=2u_{n-1}}$. \\ 
		Do $u_k\ne 0$, $\forall k\in \mathbb{N}^*$ nên nhân các  đẳng thức theo vế ta được
		$$u_1 \cdot u_2 \cdot \cdots \cdot u_n =2^{n} \cdot u_1 \cdot u_2  \cdot \cdots \cdot u_{n-1} \quad \text{hay} \quad u_n=2^n.$$
	} 
\end{ex}

% nguồn: Ok_Toan11_TracNghiem_Lan2
\begin{ex}%[1D2H1-2]%[ma:c88a9866]
	Cho dãy số có các số hạng đầu là $-1$, $1$, $-1$,  $1$, $-1$, $\ldots$ Số hạng tổng quát của dãy số này có dạng
	 \choice
		 {$u_n=1$}
		 {$u_n=-1$}
		{\True $u_n=(-1)^n$} 
		 {$u_n=(-1)^{n+1}$}
	\loigiai{ 
		Ta có  các số hạng đầu của dãy là $(-1)^1$, $(-1)^2$, $(-1)^3$, $(-1)^4$, $(-1)^5$, $\ldots$\\
		Suy  ra $u_n=(-1)^n$.
	} 
\end{ex}

% nguồn: Ok_Toan11_TracNghiem_Lan2
\begin{ex}%[1D2H1-2]%[ma:c446dddb]
	Cho dãy số có các số hạng đầu là $9$, $99$, $999$, $9999$, $\ldots$ Số hạng tổng quát  của dãy số này là
	 \choice
		{$u_n=\dfrac{n}{n+1}$} 
		{\True $u_n=10^n-1$} 
		 {$u_n=9^n$}
		 {$u_n=9n$}
	 \loigiai{
		Ta có $u_1=10^1-1$; $u_2=10^2-1$; $u_3=10^3-1$;  $u_4=10^4-1$.\\
		Vậy số  hạng tổng quát là $u_n=10^n-1$.
	 }
\end{ex}

% nguồn: Ok_Toan11_TracNghiem_Lan2
\begin{ex}%[1D2H1-2]%[ma:370ed71e]
	Cho dãy số có các số hạng đầu là $0$, $\dfrac{1}{2}$, $\dfrac{2}{3}$, $\dfrac{3}{4}$, $\dfrac{4}{5}$, $\ldots$ Số hạng tổng quát  của dãy số này là
	\choice 
		{$u_n=\dfrac{n+1}{n}$} 
		 {\True $u_n=\dfrac{n}{n+1}$}
		 {$u_n=\dfrac{n-1}{n}$}
		 {$u_n=\dfrac{n^2-n}{n+1}$}
	 \loigiai{
		 Ta có
		\begin{align*}
			0 &= \dfrac{0}{0+1}, \\ 
			\dfrac{1}{2} &= \dfrac{1}{1+1}, \\ 
			 \dfrac{2}{3} &= \dfrac{2}{2+1}, \\
			\dfrac{3}{4} &= \dfrac{3}{3+1},  \\
			\dfrac{4}{5} &=  \dfrac{4}{4+1}.
		\end{align*}
		Suy ra  $u_n=\dfrac{n}{n+1}$.
	 }
\end{ex}

% nguồn: Ok_Toan11_TracNghiem_Lan2
\begin{ex}%[1D2H1-2]%[ma:08f3a8ed]
	Cho dãy số  $(u_n)$ xác định bởi $\heva{&u_1=1 \\
	&u_{n+1}=3u_n, \ (n\ge 1)} $. Tìm công  thức số hạng tổng quát của dãy số trên.
	\choice 
		{$u_n=3^n$} 
		{\True  $u_n=3^{n-1}$}
		 {$u_n=3^{n+1}-2$}
		 {$u_n=3^n-2$}
	 \loigiai{
		Ta  có $u_1 = 1 = 3^0$, $u_2 = 3^1$, $u_3 = 3^2$, $\ldots$\\
		Dự đoán $u_n=3^{n-1}$,  $n \in \mathbb{N}^*$.\\
		Ta dễ dàng chứng minh được công  thức này bằng quy nạp.
	 }
\end{ex}

% nguồn: Ok_Toan11_TracNghiem_Lan2
\begin{ex}%[1D2H1-2]%[ma:a5702bec]
	Cho  dãy số $3$, $6$, $9$, $12$, $15$, $\ldots$ Số hạng tổng quát của dãy số này là
	\choice 
		 {$u_n = 3(n - 1)$}
		{\True $u_n =  3n$}
		{$u_n  = 3 + n$}
		{$u_n = 3n +  1$}
	\loigiai{ 
		 Ta có
		\begin{align*}
			u_1 &= 3 = 3 \cdot 1,  \\
			u_2 &=  6 = 3 \cdot 2, \\
			u_3  &= 9 = 3 \cdot 3, \\
			u_4 &= 12 = 3 \cdot 4,  \\
			u_5 &= 15  = 3 \cdot 5, \ldots
		\end{align*}
		Suy  ra $u_n = 3n$.
	 }
\end{ex}

% nguồn: Ok_Toan11_TracNghiem_Lan2
\begin{ex}%[1D2H1-2]%[ma:4c3e65d2]
	Cho dãy số  $\dfrac{1}{2}$, $\dfrac{3}{5}$, $\dfrac{2}{3}$, $\dfrac{5}{7}$, $\ldots$ Công thức tổng quát $u_n$ nào là của dãy số đã cho?
	\choice 
		{$u_n=\dfrac{n}{n+1}$, $\forall n \in \mathbb{N}^*$} 
		{$u_n=\dfrac{n}{2^n}$,  $\forall n \in \mathbb{N}^*$}
		{\True $u_n=\dfrac{n+1}{n+3}$, $\forall n  \in \mathbb{N}^*$}
		{$u_n=\dfrac{2n}{2n+1}$,  $\forall n \in \mathbb{N}^*$}
	 \loigiai{
		Viết lại dãy  số $\dfrac{2}{4}$, $\dfrac{3}{5}$, $\dfrac{4}{6}$, $\dfrac{5}{7}$, $\ldots$\\
		Suy ra $u_n=\dfrac{n+1}{n+3}$, $\forall n \in  \mathbb{N}^*$.
	} 
\end{ex}

% nguồn: Ok_Toan11_TracNghiem_Lan2
\begin{ex}%[1D2H1-2]%[ma:a7399c6c]
	Cho dãy số $(u_n)$ với $\heva{&u_1=1 \\ &u_{n+1}=u_n+(-1)^{2n+1}}$. Số hạng tổng quát  $u_n$ của dãy số là số hạng nào dưới đây?
	 \choice
		{\True  $u_n=2-n$}
		{$u_n=2+nn$} 
		{$u_n=1-n$} 
		{$u_n=-n$} 
	 \loigiai{
		Ta có $u_{n+1}=u_n+(-1)^{2n+1}=u_n-1$.  \\
		Ta có $u_1=1=2-1$,  $u_2=0=2-2$, $u_3=-1=2-3$, $\ldots$. \\
		Vậy $u_n=2-n$ với  mọi $n \in \mathbb{N}^*$.
	} 
\end{ex}

% nguồn: Ok_Toan11_TracNghiem_Lan2
\begin{ex}%[1D2H1-2]%[ma:a42830b2]
	Cho dãy số có các số  hạng đầu là $-2$, $0$, $2$, $4$, $6$, $\ldots$ Số hạng tổng quát của dãy số này có dạng
	\choice 
		{$u_n = -2n$} 
		{$u_n = -2 + n$} 
		{$u_n = -2(n+1)$} 
		{\True  $u_n = -2 + 2(n-1)$}
	\loigiai{ 
		Dãy số là dãy cách  đều với khoảng cách là $2$ và số hạng đầu tiên là $-2$. Do đó, số hạng tổng quát là $u_n = -2 + 2(n-1)$.
	 }
\end{ex}

% nguồn: Ok_Toan11_TracNghiem_Lan2
\begin{ex}%[1D2H1-2]%[ma:2ae594a5]
	Cho dãy số $(u_n)$ với $\heva{&u_1=-2 \\ &u_{n+1}=-2-\dfrac{1}{u_n}, \; n \in \mathbb{N}^*}$. 
	 Công thức số hạng tổng quát của dãy số này là
	\choice 
		{$u_n=-\dfrac{n-1}{n}$} 
		{$u_n=\dfrac{n+1}{n}$} 
		{\True  $u_n=-\dfrac{n+1}{n}$}
		 {$u_n=-\dfrac{n}{n+1}$}
	\loigiai{ 
		 Ta có
		$$u_1=-2=-\dfrac{2}{1}, \; u_2=-2+\dfrac{1}{2}=-\dfrac{3}{2}, \; u_3=-2+\dfrac{2}{3}=-\dfrac{4}{3}, \;  u_4=-2+\dfrac{3}{4}=\dfrac{5}{4}, \; \ldots.$$
		Như vậy  $u_n=-\dfrac{n+1}{n}$.
	 }
\end{ex}

% nguồn: Ok_Toan11_TracNghiem_Lan2
\begin{ex}%[1D2H1-2]%[ma:b9b42944]
	Cho dãy số $(u_n)$ có các số hạng đầu là $\dfrac{1}{2}$, $\dfrac{1}{2^2}$, $\dfrac{1}{2^3}$, $\dfrac{1}{2^4}$, $\dfrac{1}{2^5}$, $\ldots$ Số  hạng tổng quát của dãy số này là
	\choice 
		{$u_n = \dfrac{1}{2}  \cdot \dfrac{1}{2^{n+1}}$}
		 {$u_n = \dfrac{1}{2^{n+1}}$}
		 {\True $u_n = \dfrac{1}{2^n}$}
		{$u_n  = \dfrac{1}{2^{n-1}}$}
	\loigiai{ 
		 Ta có
		\begin{align*}
			u_1 &= \dfrac{1}{2^1},\\ 
			u_2 &= \dfrac{1}{2^2},\\ 
			u_3  &= \dfrac{1}{2^3},\\
			u_4  &= \dfrac{1}{2^4},\\
			u_5 &= \dfrac{1}{2^5}, \ldots 
		\end{align*}
		 Suy ra dãy $u_n = \dfrac{1}{2^n}$.
	 }
\end{ex}

% nguồn: Ok_Toan11_TracNghiem_Lan2
\begin{ex}%[1D2H1-2]%[ma:702c6fc6]
	Cho $(u_n)$ là dãy các  số tự nhiên chia hết cho $5$ viết theo thứ tự tăng dần và $u_1=5$. Dạng khai triển của dãy số $(u_n)$ là
	\choice 
		 {$0$, $5$, $10$, $\ldots$, $5n$, $\ldots$}
		{\True $5$, $10$, $15$, $\ldots$, $5n$, $\ldots$} 
		 {$5$, $15$, $25$, $\ldots$, $5n$, $\ldots$}
		{$1$, $5$,  $10$, $\ldots$, $5n$, $\ldots$}
	\loigiai{ 
		Dãy số với số hạng đầu $u_1=5$ và chia hết cho $5$ nên dạng khai triển phải là  $5$, $10$, $15$, $\ldots$, $5n$, $\ldots$
	 }
\end{ex}

% nguồn: Ok_Toan11_TracNghiem_Lan2
\begin{ex}%[1D2H1-2]%[ma:84babeca]
	Cho dãy số có các số hạng đầu là $\dfrac{1}{3}$, $\dfrac{1}{3^2}$, $\dfrac{1}{3^3}$, $\dfrac{1}{3^4}$, $\dfrac{1}{3^5}$, $\ldots$ Số hạng tổng quát của dãy  số này là
	\choice 
		{$u_n  = \dfrac{1}{3} \cdot \dfrac{1}{3^{n+1}}$}
		{$u_n =  \dfrac{1}{3^{n+1}}$}
		 {\True $u_n = \dfrac{1}{3^n}$}
		{$u_n = \dfrac{1}{3^{n-1}}$} 
	 \loigiai{
		 Ta có
		\begin{align*}
			 u_1 &= \dfrac{1}{3^1},\\
			u_2 &=  \dfrac{1}{3^2},\\
			 u_3 &= \dfrac{1}{3^3},\\
			u_4 &= \dfrac{1}{3^4},\\ 
			u_5 &= \dfrac{1}{3^5}, \ldots 
		\end{align*}
		Suy ra  dãy $u_n = \dfrac{1}{3^n}$.
	} 
\end{ex}

% nguồn: Ok_Toan11_TracNghiem_Lan2
\begin{ex}%[1D2H1-2]%[ma:bf4daecc]
	Cho dãy số $(u_n)$ với $\heva{&u_1=5 \\ &u_{n+1}=u_n+n, \; n \ge 1}$. Số hạng tổng quát $u_n$ của  dãy số là
	\choice 
		{$u_n=\dfrac{(n-1)  \cdot n}{2}$}
		 {\True $u_n=5+\dfrac{(n-1) \cdot n}{2}$}
		{$u_n=5+\dfrac{(n+1)  \cdot n}{2}$}
		 {$u_n=5+\dfrac{(n+1) \cdot (n+2)}{2}$}
	\loigiai{ 
		Theo công thức truy hồi:  $u_{n+1}-u_n=n$. Khi đó
		$$u_2-u_1=1,  \; u_3-u_2=2, \; \ldots, \; u_n-u_{n-1}=n-1.$$
		Cộng  theo vế các đẳng thức trên ta có
		$$u_n-u_1=1+2+ \ldots  +(n-1)=\dfrac{(n-1) \cdot n}{2}.$$
		Vậy $u_n=5+\dfrac{(n-1) \cdot  n}{2}$.
	 }
\end{ex}

% nguồn: Ok_Toan11_TracNghiem_Lan2
\begin{ex}%[1D2H1-2]%[ma:12f698f0]
	Cho dãy số có các số hạng đầu  là $0{,}1$; $0{,}01$; $0{,}001$; $0{,}0001$; $\ldots$ Số hạng tổng quát của dãy số này có dạng
	 \choice
		 {$u_n = \dfrac{1}{10^{n-1}}$}
		{$u_n = \dfrac{1}{10^{n+2}}$} 
		 {\True $u_n = \dfrac{1}{10^n}$}
		{$u_n  = \dfrac{1}{10^{n+1}}$}
	\loigiai{ 
		Ta  có
		\begin{align*}
			u_1 &=  0,1 = \dfrac{1}{10}, \\
			u_2 &=  0,01 = \dfrac{1}{10^2}, \\
			 u_3 &= 0,001 = \dfrac{1}{10^3}, \\
			u_4&= 0,0001 =  \dfrac{1}{10^4}.
		\end{align*}
		Dự đoán  $u_n = \dfrac{1}{10^n}$.
	} 
\end{ex}

% nguồn: Ok_Toan11_TracNghiem_Lan2
\begin{ex}%[1D2H1-2]%[ma:2c2f1920]
	Cho dãy số $(u_n)$ với $\heva{&u_1=-1 \\ &u_{n+1}=\dfrac{u_n}{2}}$. Công thức  số hạng tổng quát của dãy số này là
	\choice 
		{$u_n=(-1)  \cdot \left(\dfrac{1}{2}\right)^n$}
		{$u_n=(-1) \cdot \left(\dfrac{1}{2}\right)^{n+1}$} 
		 {$u_n=\left(\dfrac{1}{2}\right)^{n-1}$}
		{\True $u_n=(-1)  \cdot \left(\dfrac{1}{2}\right)^{n-1}$}
	 \loigiai{
		 Ta có:
		$\heva{&u_1=-1 \\ 
		&u_2=\dfrac{u_1}{2}  \\
		&u_3=\dfrac{u_2}{2}  \\
		&\vdots \\ 
		 &u_n=\dfrac{u_{n-1}}{2}}$. \\
		Do $u_k\ne 0$, $\forall k\in \mathbb{N}^*$ nên nhân các đẳng thức theo  vế ta được
		$$u_1 \cdot u_2 \cdot \cdots u_n=(-1) \cdot \dfrac{u_1 \ldots u_2 \cdot \cdots \cdot u_{n-1}}{2^{n-1}} \quad \text{hay} \quad  u_n=(-1) \cdot \left(\dfrac{1}{2}\right)^{n-1}.$$
	 }
\end{ex}

% nguồn: Ok_Toan11_TracNghiem_Lan2
\begin{ex}%[1D2H1-2]%[ma:132ced1a]
	Cho dãy số có các số hạng đầu là $8$, $15$, $22$, $29$, $36$, $\ldots$ Số hạng tổng quát của dãy  số này là
	\choice 
		{$u_n=7n+7$} 
		{$u_n=7n$} 
		 {\True $u_n=7n+1$}
		{$u_n$ không viết  được dưới dạng công thức}
	 \loigiai{
		Ta có 
		\begin{align*}
			8 &= 7 \cdot 1 + 1, \\ 
			15  &= 7 \cdot 2 + 1, \\
			22 &= 7 \cdot  3 + 1, \\
			29 &= 7 \cdot 4 +  1, \\
			36 &= 7 \cdot 5  + 1.
		\end{align*}
		Suy ra số hạng tổng quát $u_n=7n+1$. 
	} 
\end{ex}

% nguồn: Ok_Toan11_TracNghiem_Lan2
\begin{ex}%[1D2H1-2]%[ma:3b142039]
	Cho dãy số $\left(u_n \right)$ với $u_n=\dfrac{n^2+3n+7}{n+1}$, $n\in \mathbb{N}^{*}$. Hỏi dãy số trên có bao  nhiêu số hạng nhận giá trị nguyên?
	 \choice
		{$2$} 
		 {$4$}
		{\True $1$} 
		 {Không có}
	\loigiai{ 
		 Ta có $u_n=\dfrac{n^2+3n+7}{n+1}=n+2+\dfrac{5}{n+1} $. \\
		Để $u_n$ nhận giá trị nguyên thì $\dfrac{5}{n+1} \left(n\in \mathbb{N}^{*} \right)$ là  số nguyên hay $n=4$.\\
		Vậy dãy số $\left(u_n \right)$ chỉ có một số hạng nhận giá trị  nguyên.
	} 
\end{ex}

% nguồn: Ok_Toan11_TracNghiem_Lan2
\begin{ex}%[1D2H1-2]%[ma:86c5530b]
	Cho dãy số  $(u_n)$ biết $u_n=3^n$. Tìm số hạng $u_{2n-1}$.
	 \choice
		{$u_{n-1}=5^{n-1}$} 
		{\True $u_{2n-1}=3^n \cdot  3^{n-1}$}
		 {$u_{2n-1}= 3^{2n}-1$}
		 {$u_{2n-1}=3^{2(n-1)}$}
	\loigiai{Số  hạng $u_{2n-1}$ là $u_{2n-1}=3^{2n-1}=3^{n+(n-1)}=3^{n} \cdot 3^{n-1}$.
	 }
\end{ex}

% nguồn: Ok_Toan11_TracNghiem_Lan2
\begin{ex}%[1D2H1-2]%[ma:aba698a3]
	Cho  dãy số \((u_n)\) với $\heva{&u_1=\dfrac{1}{2} \\ &u_{n+1}=u_n-2}$. Công thức số hạng tổng quát của dãy số này là
	 \choice
		{$u_n=\dfrac{1}{2}+2(n-1)$} 
		{\True $u_n=\dfrac{1}{2}-2(n-1)$} 
		 {$u_n=\dfrac{1}{2}-2n$}
		{$u_n=\dfrac{1}{2}+2n$} 
	\loigiai{ 
		Ta  có
		$\heva{&u_1=\dfrac{1}{2} \\ 
		&u_2=u_1-2 \\ 
		 &u_3=u_2-2 \\
		&\vdots \\ 
		&u_n=u_{n-1}-2.}$\\ 
		Cộng các  đẳng thức trên theo vế ta được
		\[u_1+u_2+u_3+\cdots+u_n=\dfrac{1}{2}+u_1-2+u_2-2+\cdots+u_{n-1}-2.\] 
		Thu gọn lại ta được số hạng  tổng quát là $u_n=\dfrac{1}{2}-2(n-1)$.
	} 
\end{ex}

% nguồn: Ok_Toan11_TracNghiem_Lan2
\begin{ex}%[1D2H1-2]%[ma:97bf8a4f]
	 Cho dãy số $(u_n)$, với $u_n=(n-1)^{2023}$. Tìm số hạng $u_{n+1}$
	 \choice
		 {$u_{n+1}=(n+1)^{2023}$}
		{$u_{n+1}=(n+1)^{2022}$} 
		 {$u_{n+1}=(n-1)^{2023}$}
		{\True  $u_{n+1}=n^{2023}$}
	\loigiai{ 
		Ta có $u_{n+1}=[(n+1)-1]^{2023}=n^{2023}$. 
	 }
\end{ex}

% nguồn: Ok_Toan11_TracNghiem_Lan2
\begin{ex}%[1D2H1-2]%[ma:e2e129a4]
	Cho dãy số $(u_n)$ xác định bởi $\heva{&u_1=1 \\  &u_{n+1}=u_n+2, \; n \ge 1}$. Số hạng tổng quát của dãy số trên là
	\choice 
		 {\True $u_n=2n-1$}
		 {$u_n=3n-2$}
		 {$u_n=4n-3$}
		{$u_n=8n-7$} 
	 \loigiai{
		 Ta có
		\begin{eqnarray*}
			&&u_1=1=2\cdot  1-1;\\
			&&u_2=3=2\cdot  2-1;\\
			 &&u_3=5=2\cdot 3-1
		\end{eqnarray*}
		Từ đó suy ra số hạng tổng quát của  dãy số là $u_n=2n-1$.
	 }
\end{ex}

% nguồn: Ok_Toan11_TracNghiem_Lan2
\begin{ex}%[1D2H1-2]%[ma:1eb104a7]
	Cho dãy  số có các số hạng đầu là
	$\dfrac{1}{3}$; $\dfrac{1}{3^2}$; $\dfrac{1}{3^3}$; $\dfrac{1}{3^4}$; $\ldots$. Số hạng tổng quát của  dãy số này là
	 \choice
		{$u_n=\dfrac{1}{3} \cdot  \dfrac{1}{3^{n+1}}$}
		{$u_n=\dfrac{1}{3^{n+1}}$} 
		 {\True $u_n=\dfrac{1}{3^n}$}
		 {$u_n=\dfrac{1}{3^{n-1}}$}
	 \loigiai{
		Ta  có $u_1=\dfrac{1}{3}$; $u_2=\dfrac{1}{3^2}$; $u_3=\dfrac{1}{3^3}$; $u_4=\dfrac{1}{3^4}$; $\ldots$.\\
		Suy ra  được công thức tổng quát $u_n=\dfrac{1}{3^n}$.
	} 
\end{ex}

% nguồn: Ok_Toan11_TracNghiem_Lan2
\begin{ex}%[1D2H1-2]%[ma:99abc081]
	Cho dãy số có các số hạng đầu là: $-2$; $0$; $2$; $4$; $6$; $\ldots$. Số hạng tổng quát  của dãy số này là
	\choice 
		 {$u_n=-2n$}
		{$u_n=(-2)+n$} 
		 {$u_n=-2(n+1)$}
		{\True  $u_n=-2+2(n-1)$}
	 \loigiai{
		 Thay giá trị $n=1$ vào từng phương án thì chỉ có $u_n=-2+(n-1)$ thỏa mãn.
	} 
\end{ex}

% nguồn: Ok_Toan11_TracNghiem_Lan2
\begin{ex}%[1D2H1-2]%[ma:bb63f420]
	Cho dãy số viết dưới  dạng khai triển là $1$, $4$, $9$, $16$, $25$. Trong các công thức sau, công thức nào là công thức tổng quát của dãy số trên.
	\choice 
		{$u_n =  3n - 2$}
		{$u_n =  n + 3$}
		{\True $u_n = n^2$} 
		{$u_n = 2n^2 - 1$} 
	 \loigiai{
		Ta có $1 = 1^2$, $4 = 2^2$, $9 = 3^2$,  $16 = 4^2$, $25 = 5^2$.\\
		Vậy dãy số có công thức  tổng quát $u_n = n^2$.
	 }
\end{ex}

% nguồn: Ok_Toan11_TracNghiem_Lan2
\begin{ex}%[1D2H1-2]%[ma:2dda5dc0]
	Cho  dãy số $(u_n)$ với $\heva{&u_1=2 \\ &u_{n+1}-u_n=2n-1}$. Số hạng tổng quát $u_n$ của dãy số là số hạng nào dưới đây?
	\choice 
		 {\True $u_n=2+(n-1)^2$}
		 {$u_n=2+n^2$}
		{$u_n=2+(n+1)^2$} 
		 {$u_n=2-(n-1)^2$}
	 \loigiai{
		Ta  có
		$\heva{&u_1=2  \\
		 &u_2=u_1+1 \\
		&u_3=u_2+3 \\ 
		 &\vdots \\
		&u_n=u_{n-1}+(2n-3)}$\\ 
		Cộng các đẳng thức trên theo vế ta  có
		\[u_1+u_2+u_3+\cdots+u_n=2+u_1+1+u_2+3+u_3+5+\cdots+u_{n-1}+(2n-3)\] 
		Thu  gọn lại ta được
		$$u_n=2+1+3+5+  \ldots +(2n-3)=2+(n-1)^2.$$
	} 
\end{ex}

% nguồn: Ok_Toan11_TracNghiem_Lan2
\begin{ex}%[1D2H1-2]%[ma:24499c76]
	Cho dãy số có các số hạng đầu là $5$, $10$, $15$, $20$, $25$,  $\ldots$ Số hạng tổng quát của dãy số này là
	 \choice
		{$u_n=5(n-1)$} 
		{\True  $u_n=5n$}
		 {$u_n=5+n$}
		{$u_n=5n+1$} 
	 \loigiai{
		Ta có 
		\begin{align*}
			5 &= 5 \cdot 1, \\ 
			10 &= 5 \cdot 2, \\ 
			15 &= 5 \cdot  3, \\
			20 &=  5 \cdot 4, \\
			25 &= 5 \cdot  5.
		\end{align*}
		Suy ra số hạng tổng quát  $u_n=5n$.
	 }
\end{ex}

% nguồn: Ok_Toan11_TracNghiem_Lan2
\begin{ex}%[1D2H1-2]%[ma:793eadf2]
	Cho dãy số $(u_n)$ với $\heva{&u_1=\dfrac{1}{2} \\ &u_{n+1}=2u_n, \; n \ge  2}$. Công thức số hạng tổng quát của dãy số này là
	\choice 
		 {$u_n=-2^{n-1}$}
		{$u_n=-\dfrac{1}{2^{n-1}}$} 
		{$u_n=-\dfrac{1}{2^n}$} 
		 {\True $u_n=2^{n-2}$}
	\loigiai{ 
		Ta có 
		$\heva{&u_1=\dfrac{1}{2}  \\
		&u_2=2u_1 \\ 
		&u_3=2u_2 \\ 
		&\vdots  \\
		&u_n=2u_{n-1}}$. \\ 
		Do $u_k\ne 0$, $\forall k\in \mathbb{N}^*$  nên nhân các đẳng thức theo vế ta được
		$$u_1 \cdot u_2 \cdot  \cdots \cdot u_n =\dfrac{1}{2} \cdot u_1 \cdot u_2 \cdot \cdots \cdot u_{n-1} \quad \text{hay} \quad u_n=2^{n-2}.$$
	} 
\end{ex}

% nguồn: Ok_Toan11_TracNghiem_Lan2
\begin{ex}%[1D2H1-2]%[ma:6385618a]
	Cho dãy số $(u_n)$  với $u_n=5^{n+1}$. Tìm số hạng $u_{n-1}$.
	\choice 
		 {$u_{n-1}=5^{n-1}$}
		 {\True $u_{n-1}=5^n$}
		{$u_{n-1}=5  \cdot 5^{n+1}$}
		{$u_{n-1}=5 \cdot  5^{n-1}$}
	\loigiai{Số hạng $u_{n-1}$ là $u_{n-1}=5^{(n-1)+1}=5^n$. 
	 }
\end{ex}

% nguồn: Ok_Toan11_TracNghiem_Lan2
\begin{ex}%[1D2H1-2]%[ma:cdff434f]
	Cho  dãy số $(u_n)$ biết $u_n=2^n$. Tìm số hạng $u_{n+1}$.
	 \choice
		{$u_{n+1}=2(n+1)$} 
		 {$u_{n+1}=2^n+1$}
		{\True $u_{n+1}=2^n \cdot  2$}
		{$u_{n+1}=2^n+2$} 
	\loigiai{Ta có $u_{n+1}=2^{n+1}=2^n \cdot  2$.
	 }
\end{ex}

% nguồn: Ok_Toan11_TracNghiem_Lan2
\begin{ex}%[1D2H1-2]%[ma:8d2ef806]
	 Cho dãy số $\left(u_n\right)$ với $\heva{&u_1=1 \\ &u_{n+1}=u_n+(-1)^{2n}}$.
	Số hạng tổng quát $u_n$  của dãy số là số hạng nào dưới đây?
	 \choice
		{$u_n=1+n$} 
		{$u_n=1-n$} 
		 {$u_n=1+(-1)^{2n}$}
		{\True  $u_n=n$}
	\loigiai{ 
		 Ta có $u_{n+1}=u_n+(-1)^{2n}=u_n+1$. \\
		Khi  đó, $u_1=1$, $u_2=u_1+1=2$, $u_3=u_2+1=3$, $\ldots$. \\
		 Vậy $u_n=n$ với mọi $n \in \mathbb{N}^*$.
	 }
\end{ex}

% nguồn: Ok_TracNghiemLop11_Lan1
\begin{ex}%[1D2H1-2]%[ma:cf2a831b]
	Cho dãy số $(u_n)$, biết $\heva{&u_1=-1\\&u_{n+1}=2u_n+1}$ ($n\in\mathbb{N}^{*}$). Số hạng thứ hai của dãy  số bằng
	\choice 
		{$2$} 
		 {$3$}
		 {$1$}
		 {\True $-1$}
	 \loigiai{
		Ta có  $u_2=2\cdot u_1+1=2\cdot (-1)+1=-1$.
	} 
\end{ex}

% nguồn: Ok_Toan11_TracNghiem_Lan2
\begin{ex}%[1D2V1-2]%[ma:ab2b0019]
	Cho dãy số $(u_n)$ xác định bởi $\heva{&u_1=-2 \\  &u_{n+1}=-2-\dfrac{1}{u_n}, \; n \in \mathbb{N}^*}$. Công thức số hạng tổng quát của dãy số $(u_n)$ là
	\choice 
		 {$u_n=-\dfrac{3n-1}{n}$}
		 {$u_n=-\dfrac{n}{n+1}$}
		{$u_n=\dfrac{n+1}{n}$} 
		 {\True $u_n=-\dfrac{n+1}{n}$}
	 \loigiai{
		Ta  có
		\begin{eqnarray*}
			&&u_{n+1}=-2-\dfrac{1}{u_n}  \\
			 &\Leftrightarrow& u_{n+1}+1 =-1-\dfrac{1}{u_n} \\
			&\Leftrightarrow&  u_{n+1}+1=-\dfrac{u_n+1}{u_n} \\
			&\Leftrightarrow& \dfrac{1}{u_{n+1}}=\dfrac{u_n}{u_n+1}  \\
			&\Leftrightarrow& \dfrac{1}{u_{n+1}}=\dfrac{1}{u_n+1}-1  \\
			&\Leftrightarrow& \dfrac{1}{u_n+1}-\dfrac{1}{u_{n+1}+1}=1. 
		\end{eqnarray*}
		 Đặt $v_n=\dfrac{1}{u_n+1}$. \\
		Khi đó $v_{n}-v_{n+1}=1$ hay  $v_{n+1}-v_n=-1$. \\
		Ta có $\heva{&v_2-v_1=-1 \\ &v_3-v_2=-1  \\ & \vdots \\ &v_n-v_{n-1}=-1}$. \\
		Cộng các đẳng thức trên theo vế ta  có
		\[v_2-v_1+v_3-v_2+\cdots+v_n-v_{n-1}=-1+(-1)+\cdots+(-1).\] 
		Thu gọn lại ta có $v_n-v_1=(n-1) \cdot (-1)$ hay  $v_n=v_1-n+1=\dfrac{1}{-2+1}+n+1=-n$. \\
		Do  đó $-n=\dfrac{1}{u_n+1}$ hay $u_n=-\dfrac{1}{n}-1=-\dfrac{n+1}{n}$.
	} 
\end{ex}

% nguồn: Ok_Toan11_TracNghiem_Lan2
\begin{ex}%[1D2V1-2]%[ma:2a7f411b]
	Cho dãy số $(u_n)$  với $\heva{&u_1=-1\\
	&u_{n+1}=\dfrac{u_n}{2} 
	\quad (n \ge  1)}$. Công thức số hạng tổng quát của dãy số này là
	\choice 
		 {$u_n=(-1) \cdot\left(\dfrac{1}{2}\right)^n$}
		 {$u_n=(-1) \cdot\left(\dfrac{1}{2}\right)^{n+1}$}
		{$u_n=\left(\dfrac{1}{2}\right)^{n-1}$} 
		 {\True $u_n=(-1)\cdot\left(\dfrac{1}{2}\right)^{n-1}$}
	 \loigiai{
		Ta có  $\heva{&u_1=-1\\
		&u_2=\dfrac{u_1}{2} \\ 
		&u_3=\dfrac{u_2}{2} \\ 
		 &\ldots \\
		 &u_n=\dfrac{u_{n-1}}{2}.
		 }$ \\
		Nhân hai vế ta được $u_1 \cdot u_2  \cdot u_3 \ldots u_n=(-1) \cdot \frac{u_1 \cdot u_2 \cdot u_3 \ldots u_{n-1}}{\underbrace{2 \cdot 2 \cdot 2}_{n-1 \text { lần }}} \Leftrightarrow u_n=(-1) \cdot \dfrac{1}{2^{n-1}}=(-1) \cdot\left(\dfrac{1}{2}\right)^{n-1}$.
	 }
\end{ex}

% nguồn: Ok_Toan11_TracNghiem_Lan2
\begin{ex}%[1D2V1-2]%[ma:e4e2b3dd]
	Cho  dãy số $(u_n)\colon \heva{&u_1=1\\
	&u_{n+1}=u_n+2n+1, n \geq 1}$. Giá  trị của $n$ để $-u_n+2\,017n+2\,018=0$ là\\
	 \choice
		 {Không có $n$}
		{$1\,009$} 
		{\True  $2\,018$}
		 {$2\,017$}
	 \loigiai{Ta có:
		\[ 
		\begin{aligned}
			&  u_1=1\\
			 & u_2=u_1+2\cdot1+1\\
			& u_3=u_2+2\cdot2+1\\ 
			&  u_4=u_3+2\cdot3+1\\
			& \cdots  \\
			 & u_n=u_{n-1}+2(n-1)+1.
		\end{aligned}
		 \]
		Cộng vế với vế các đẳng thức trên ta được 
		\allowdisplaybreaks 
		\begin{eqnarray*}
			u_n=2(1+2+3+\cdots+n-1)+n=2\cdot  \dfrac{(n-1+1)(n-1)}{2}+n=n^2, \forall n \in \mathbb{N}^{*}.
		\end{eqnarray*}
		Do đó $-u_n+2\,017n+2\,018=0\Leftrightarrow-n^2+2\,017n+2\,018=0\Leftrightarrow\hoac{&n=-1\\ 
		 &n=2\,018.}$ \\
		Vậy  $n=2\,018$.
	 }
\end{ex}

% nguồn: Ok_Toan11_TracNghiem_Lan2
\begin{ex}%[1D2V1-2]%[ma:adfef47f]
	Cho dãy số $(u_n)$ có dạng khai triển là $1$, $\dfrac{1}{5}$, $\dfrac{1}{5^2}$, $\dfrac{1}{5^3}$,  $\dfrac{1}{5^4}$, \ldots. Số hạng tổng quát $u_n$ của dãy có thể cho bởi công thức nào sau đây?
	\choice 
		 {$u_n=\dfrac{1}{5^n}$}
		{$u_n=\dfrac{1}{5^{n+1}}$} 
		 {\True $u_n=\dfrac{1}{5^{n-1}}$}
		 {$u_n=\dfrac{1}{5^{n-2}}$}
	 \loigiai{Ta có
		\begin{itemize}
			 \item $u_1=1=\dfrac{1}{5^0}=\dfrac{1}{5^{1-1}}$,
			 \item $u_2=\dfrac{1}{5}=\dfrac{1}{5^{2-1}}$,
			 \item $u_3=\dfrac{1}{5^2}=\dfrac{1}{5^{3-1}}$,
			\item $u_4=\dfrac{1}{5^3}=\dfrac{1}{5^{4-1}} \ldots$ 
		\end{itemize}
		Từ dây, ta có thể dự đoán $u_n=\dfrac{1}{5^{n-1}}$. 
	} 
\end{ex}

% nguồn: Ok_Toan11_TracNghiem_Lan2
\begin{ex}%[1D2V1-2]%[ma:04ec57ef]
	Cho dãy số  $(u_n)$ với $\heva{&u_1 = \dfrac{1}{2} \\
	&u_{n+1} = u_n - 2, \ n \in \mathbb{N}^*.}$ Công thức số hạng tổng quát của  dãy số này là:
	 \choice
		{$u_n  = \dfrac{1}{2} + 2(n-1)$}
		 {\True $u_n = \dfrac{1}{2} - 2(n-1)$}
		{$u_n =  \dfrac{1}{2} - 2n$}
		{$u_n = \dfrac{1}{2} +  2n$}
	 \loigiai{
		 Theo công thức truy hồi, ta có $u_{n+1} - u_n = - 2$. Khi đó
		\begin{align*}
			u_1 &=  \dfrac{1}{2}, \\
			u_2 -  u_1 &= -2, \\
			u_3  - u_2 &= -2, \\
			 &\ldots \\
			 u_n - u_{n-1} &= -2.
		\end{align*}
		Cộng vế theo vế các đẳng  thức trên, ta được
		\[ 
		u_n = \dfrac{1}{2} +  (-2)\cdot(n-1) = \dfrac{1}{2} - 2(n-1).
		\] 
	} 
\end{ex}

% nguồn: Ok_Toan11_TracNghiem_Lan2
\begin{ex}%[1D2V1-2]%[ma:4b572883]
	Cho dãy số có các số hạng đầu là $8$, $13$, $18$, $23$, $28$, $\ldots$ Số hạng tổng quát của dãy số này là 
	\choice 
		{$u_n  = 5n + 3$}
		{\True $u_n =  5n - 1$}
		 {$u_n = 5n + 2$}
		{$u_n  = 5n - 3$}
	\loigiai{ 
		\textbf{Cách  1:}\\
		Ta  có
		\begin{align*}
			u_1 &=  8 = 5 \cdot 1 - 1, \\
			u_2 &= 13 = 5  \cdot 2 - 1, \\
			u_3 &= 18  = 5 \cdot 3 - 1, \\
			u_4 &= 23 = 5 \cdot 4 -  1, \\
			u_5 &= 28 =  5 \cdot 5 - 1.
		\end{align*}
		Suy ra số hạng tổng quát của dãy là $u_n =  5n - 1$.\\
		\textbf{Cách  2:}\\
		Dãy số đã cho là một cấp  số cộng với $u_1 = 8$ và công sai $d = 5$.\\
		Vậy số  hạng tổng quát là $u_n = 8 + 5(n - 1) = 5n - 1$.
	} 
\end{ex}

% nguồn: Ok_Toan11_TracNghiem_Lan2
\begin{ex}%[1D2V1-2]%[ma:8c402e26]
	Cho dãy số $(u_n)$, với $u_1=6$; $u_n=u_{n-1}+5, \, \forall n \geq 2$. Khi đó  $u_n$ có thể được tính theo biểu thức nào dưới đây?
	\choice 
		 {$u_n=5(n-1)$}
		 {$u_n=5^{n+1}$}
		 {$u_n=5(n+1)$}
		{\True  $u_n=5n+1$}
	 \loigiai{
		Dãy số $(u_n)$ là cấp số cộng với $u_1=6$,  công sai $d=5$ nên có số hạng tổng quát
		\allowdisplaybreaks 
		\begin{eqnarray*}
			u_n&=&u_1+(n-1) d\\ 
			&=&6+(n-1)  5=5n+1.
		\end{eqnarray*}
	} 
\end{ex}

% nguồn: Ok_Toan11_TracNghiem_Lan2
\begin{ex}%[1D2V1-2]%[ma:ce074faf]
	 Cho dãy số $(u_n)$ với $\heva{&u_1=5 \\
	&u_{n+1}=u_n+n, \ n \in  \mathbb{N}^*.}$\\
	 Số hạng tổng quát $u_n$ của dãy số này là
	 \choice
		{$u_n = \dfrac{(n-1)n}{2}$} 
		{\True $u_n = 5 +  \dfrac{(n-1)n}{2}$}
		{$u_n  = 5 + \dfrac{(n+1)n}{2}$}
		{$u_n = 5 +  \dfrac{(n+1)(n+2)}{2}$}
	 \loigiai{
		Theo công thức truy hồi, ta có $u_{n+1} - u_n = n$.  Khi đó
		\begin{align*}
			u_1 &= 5,  \\
			u_2 -  u_1 &= 1, \\
			u_3 -  u_2 &= 2, \\
			 &\ldots\\
			u_n - u_{n-1}  &= n-1.
		\end{align*}
		Cộng vế theo vế các đẳng thức trên, ta được 
		 \[
		u_n = 5  + (1 + 2 + 3 + \ldots + (n-1)) = 5 + \dfrac{(n-1)n}{2}.
		 \]
	 }
\end{ex}

% nguồn: Ok_Toan11_TracNghiem_Lan2
\begin{ex}%[1D2V1-2]%[ma:9bbb1fa6]
	Cho  dãy số $(u_n)$ với $\heva{&u_1=1\\
	&u_{n+1}=u_n+(-1)^{2n} 
	\quad (n \ge 1)}$. Công thức tổng quát $u_n$ nào dưới đây là của  dãy số đã cho?
	\choice 
		 {\True $u_n=n$}
		 {$u_n=1-n$}
		 {$u_n=1+(-1)^{2n}$}
		 {$u_n=1+n$}
	 \loigiai{Ta có $u_{n+1}=u_n+(-1)^{2n}=u_n+1\Rightarrow u_2=2$; $u_3=3$; $u_4=4$; $\ldots$\\
		Dự đoán được $u_n=n$, $\forall n \in \mathbb{N}^{*}$. \\ 
		Ta chứng minh $u_n=n$,  $\forall n \in \mathbb{N}^{*}(*)$ bằng phương pháp quy nạp.
		\begin{itemize}
			\item Với $n=1\Rightarrow u_1=1$. Vậy  $(*)$ đúng với $n=1$.
			\item Giả sử $(*)$ đúng với $n=k\left(k \in \mathbb{N}^{*}\right)$,  tức là ta có $u_k=k$.
			\item  Ta đi chứng minh (*) cũng đúng với $n=k+1$, tức là cần chứng minh $u_{k+1}=k+1$.\\
			Thật vậy, từ hệ thức xác định dãy số $(u_n)$  ta có $u_{k+1}=u_k+(-1)^{2k}=k+1$.
		\end{itemize}
		Vậy $(*)$ đúng  với mọi $n \in \mathbb{N}^{*}$.
	} 
\end{ex}

% nguồn: Ok_Toan11_TracNghiem_Lan2
\begin{ex}%[1D2V1-2]%[ma:a14f9077]
	 Gọi $S_n=\dfrac{1}{1\cdot3}+\dfrac{1}{3\cdot 5}+\dfrac{1}{5\cdot7}+\cdots+\dfrac{1}{(2n-1)(2n+1)}$với mọi $n \in \mathbb{N}^{*}$. Ta có
	 \choice
		 {$S_n=\dfrac{n-1}{2n-1}$}
		 {$S_n=\dfrac{2n}{2n+1}$}
		{\True  $S_n=\dfrac{n}{2n+1}$}
		{$S_n=\dfrac{n+1}{2n+3}$} 
	\loigiai{Ta  có:
		\allowdisplaybreaks 
		\begin{eqnarray*}
			S_n&=&\dfrac{1}{1\cdot3}+\dfrac{1}{3\cdot5}+\dfrac{1}{5\cdot7}+\cdots+\dfrac{1}{(2n-1)(2n+1)}  \\
			 &=&\dfrac{1}{2}\left(\dfrac{1}{1}-\dfrac{1}{3}+\dfrac{1}{3}-\dfrac{1}{5}+\dfrac{1}{5}-\dfrac{1}{7}+\cdots+\dfrac{1}{(2n-1)}-\dfrac{1}{(2n+1)}\right)\\
			 &=&\dfrac{1}{2}\left(\dfrac{1}{1}-\dfrac{1}{2n+1}\right)=\dfrac{n}{2n+1}.
		\end{eqnarray*}
	 }
\end{ex}

% nguồn: Ok_Toan11_TracNghiem_Lan2
\begin{ex}%[1D2V1-2]%[ma:fb741e0a]
	Cho dãy số $(u_n)$ với $\heva{&u_1=1 \\ &u_{n+1}=u_n+n^2}$. Số hạng tổng quát $u_n$ của dãy số là số hạng  nào dưới đây?
	 \choice
		{\True $u_n=1+\dfrac{n(n-1)(2n-1)}{6}$} 
		 {$u_n=1+\dfrac{n(n+1)(2n+1)}{6}$}
		{$u_n=1+\dfrac{n(n-1)(2n+2)}{6}$} 
		{$u_n=1+\dfrac{n(n+1)(2n-2)}{6}$} 
	\loigiai{ 
		 Ta có:
		$\heva{&u_1=1  \\
		&u_2=u_1+1^2  \\
		 &u_3=u_2+2^2 \\
		 &\vdots \\
		&u_n=u_{n-1}+(n-1)^2}$\\ 
		Cộng  các đẳng thức trên theo vế
		$$u_n=1 +1^2+2^2+  \cdots +(n-1)^2=1+\dfrac{n(n-1)(2n-1)}{6}.$$
	 }
\end{ex}

% nguồn: Ok_Toan11_TracNghiem_Lan2
\begin{ex}%[1D2V1-2]%[ma:1164b0c5]
	 Cho dãy số $(u_n)$ với $\heva{&u_1=\dfrac{1}{2} \\
	&u_{n+1}=2u_n \quad (n \ge 1)}$. Công thức số hạng tổng quát của dãy số này: 
	 \choice
		 {$u_n=-2^{n-1}$}
		{$u_n=\dfrac{-1}{2^{n-1}}$} 
		{$u_n=\dfrac{-1}{2^n}$} 
		{\True $u_n=2^{n-2}$} 
	\loigiai{ 
		Ta có  $\heva{&u_1=\dfrac{1}{2} \\
		&u_2=2u_1\\ 
		 &u_3=2u_2\\
		&\cdots  \\
		&u_n=2u_{n-1}.}$\\ 
		Nhân hai vế ta được $u_1\cdot u_2\cdot u_3\ldots u_n=\dfrac{1}{2} \cdot 2^{n-1} \cdot u_1\cdot u_2\ldots u_{n-1}  \Leftrightarrow u_n=2^{n-2}$.
	 }
\end{ex}

% nguồn: Ok_Toan11_TracNghiem_Lan2
\begin{ex}%[1D2V1-2]%[ma:f488758f]
	Cho  dãy số $(u_n)$ với $\heva{&u_1=2\\
	&u_{n+1}-u_n=2n-1 \quad (n  \ge 1)}$. Số hạng tổng quát $u_n$ của dãy số là số hạng nào dưới đây?
	 \choice
		{\True  $u_n=2+(n-1)^2$}
		{$u_n=2+n^2$} 
		{$u_n=2+(n+1)^2$} 
		 {$u_n=2-(n-1)^2$}
	\loigiai{ 
		Ta có: $\heva{&u_1=2\\ 
		&u_2=u_1+1\\ 
		&u_3=u_2+3\\ 
		&\ldots  \\
		 &u_n=u_{n-1}+2n-3.}$ \\
		 Cộng vế theo ta được
		$u_n=2+1+3+5+\cdots+(2n-3)=2+(n-1)^2$. 
	} 
\end{ex}

% nguồn: Ok_Toan11_TracNghiem_Lan2
\begin{ex}%[1D2V1-2]%[ma:56a67ae0]
	Cho  dãy số $(u_n)$ với $\heva{&u_1=1\\
	&u_{n+1}=u_n+n^2 \quad (n \ge 1)}$. Số hạng tổng quát $u_n$ của dãy số là số hạng  nào dưới đây?
	 \choice
		{$u_n=1+\dfrac{n(n+1)(2n+1)}{6}$} 
		 {$u_n=1+\dfrac{n(n-1)(2n+2)}{6}$}
		{\True  $u_n=1+\dfrac{n(n-1)(2n-1)}{6}$}
		{$u_n=1+\dfrac{n(n+1)(2n-2)}{6}$} 
	\loigiai{Ta có $\heva{&u_1=1\\ 
		 &u_2=u_1+1^2\\
		 &u_3=u_2+2^2\\
		&\cdots \\ 
		&u_n=u_{n-1}+(n-1)^2.}$  \\
	Cộng hai  vế ta được $u_n=1+1^2+2^2+\cdots+(n-1)^2=1+\dfrac{n(n-1)(2n-1)}{6}$.}
\end{ex}

% nguồn: Ok_Toan11_TracNghiem_Lan2
\begin{ex}%[1D2V1-2]%[ma:058d7d28]
	Cho dãy  số $(u_n)$ xác định bởi $\heva{&u_1=1 \\
	&u_{n+1}=u_n+2, \ (n \ge 1).}$\\ 
	Xác định  công thức số hạng tổng quát của dãy số này.
	 \choice
		{\True $u_n = 2n  - 1$}
		{$u_n = 3n  - 2$}
		{$u_n  = 4n - 3$}
		{$u_n  = 8n - 7$}
	\loigiai{ 
		Theo công thức truy hồi, ta có $u_{n+1} -  u_n = 2$. Khi đó:
		\begin{align*}
			u_1 &= 1,  \\
			u_2 - u_1 &= 2,  \\
			u_3 -  u_2 &= 2, \\
			&\ldots\\ 
			u_n  - u_{n-1} &= 2.
		\end{align*}
		Cộng vế theo vế  các đẳng thức trên, ta được
		 \[
		u_n = 1 + (2 + 2 + 2 + \ldots +  2) = 2n - 1.
		 \]
	} 
\end{ex}

% nguồn: Ok_Toan11_TracNghiem_Lan2
\begin{ex}%[1D2V1-2]%[ma:f9818757]
	Cho hai cấp số cộng  $u_n\colon 1$; $6$; $11$; $\ldots$ và $v_n\colon 4$; $7$; $10$; $\ldots$. Mỗi cấp số cộng có $2\,018$ số hạng. Hỏi có bao nhiêu số có mặt trong cả hai dãy số trên?
	 \choice
		{\True $403$} 
		 {$401$}
		{$402$} 
		 {$504$}
	\loigiai{ 
		\begin{itemize}
			\item Cấp số cộng $(u_n)$ có số hạng tổng quát là $u_n=1+5n-1=5n-4, \quad 1\leq n \leq 2\,018$. 
			\item Cấp số cộng  $(v_m)$ có số hạng tổng quát là $v_m=4+3m-1=3m+1, \quad 1\leq m \leq 2\,018$.
		\end{itemize}
		Một số có mặt trong cả hai dãy số trên nếu  tồn tại $m, n \in \mathbb{N}$ thỏa mãn điều kiện $\heva{&1\leq m, n \leq 2\,018\\
		 &u_m=u_n\left(^{*}\right).
		}$\\ 
		Ta  có $\left(^{*}\right) \Leftrightarrow 5n-4=3m+1\Leftrightarrow 5n-5=3m \left(^{**}\right)$. \\
		Do  $3$ và $5$ là hai số nguyên tố cùng nhau nên từ $\left(^{**}\right)$ suy ra $m\, \vdots \, 5$ tức tồn tại $k \in \mathrm{Z}$ sao cho $m=5k$. \\
		Mặt khác $1\leq m \leq 2\,018 \Leftrightarrow 1 \le 5k \le  2\,018 \Leftrightarrow \dfrac{1}{5} \le k \le \dfrac{2\,018}{5}$. \\
		Mà $k \in \mathrm{Z}$ nên $k \in \{1;2;\ldots;672\}$, khi đó ta có $m \in \{5; 10;  \ldots; 2\,015\}$.\\
		Mặc khác, thay $m=5k$  vào $(^{**})$ được
		 \allowdisplaybreaks
		\begin{eqnarray*}
			5n-5=3(5k)  \Leftrightarrow n=3k+1.
		\end{eqnarray*}
		Với $k=403$, ta được $n=3\cdot 403+1=1\,210<2\,018$. Do đó  mỗi giá trị $k$ sẽ có $1$ giá trị $n$ tương ứng. \\
		Vậy trong $2\,018$ số hạng đầu  tiên ở cả hai cấp số cộng, có $403$ số có mặt trong cả hai dãy đã cho.
	 }
\end{ex}

% nguồn: Ok_Toan11_TracNghiem_Lan2
\begin{ex}%[1D2V1-2]%[ma:4c179af6]
	Cho dãy  số $(u_n)$ xác định bởi $\heva{&u_1=-2\\
	&u_{n+1}=-2-\dfrac{1}{u_n}, \forall  n \in \mathbb{N}^{*}}$. Tìm công thức số hạng tổng quát của dãy số.
	\choice 
		 {$u_n=-\dfrac{3n-1}{n}$}
		 {$u_n=-\dfrac{n}{n+1}$}
		{$u_n=\dfrac{n+1}{n}$} 
		{\True $u_n=-\dfrac{n+1}{n}$} 
	\loigiai{ 
		Ta  có
		 \allowdisplaybreaks
		\begin{eqnarray*}
			u_{n+1}=-2-\dfrac{1}{u_n}  &\Leftrightarrow& u_{n+1} \cdot u_n=-2u_n-1\\
			&\Leftrightarrow& u_{n+1} \cdot  u_n+u_n+u_{n+1}+1=u_{n+1}-u_n \\
			&\Leftrightarrow&\left(u_{n+1}+1\right)\left(u_n+1\right)=u_{n+1}-u_n  \\
			&\Leftrightarrow&  \dfrac{u_{n+1}-u_n}{\left(u_{n+1}+1\right)\left(u_n+1\right)}=1 \\
			&\Leftrightarrow&  \dfrac{1}{u_n+1}-\dfrac{1}{u_{n+1}+1}=1.
		\end{eqnarray*}
		Đặt $v_n=\dfrac{1}{u_n+1}$. Khi đó $v_{n+1}-v_n=-1$ nên $(v_n)$ là một cấp số cộng  với công sai $d=-1$. \\
		Ta  có
		\allowdisplaybreaks 
		\begin{eqnarray*}
			 v_n=v_1+(n-1)(-1) &\Leftrightarrow& v_n=\dfrac{1}{u_1+1}-n+1=-n\\
			&\Leftrightarrow& \dfrac{1}{u_n+1}=-n  \\
			 &\Leftrightarrow& u_n+1=-\dfrac{1}{n}\\
			&\Leftrightarrow& u_n=-\dfrac{1}{n}-1=-\dfrac{n+1}{n}. 
		\end{eqnarray*}
	 }
\end{ex}

% nguồn: Ok_Toan11_TracNghiem_Lan2
\begin{ex}%[1D2V1-2]%[ma:b6ba7d7c]
	 Cho dãy số $(u_n)$ với $\heva{&u_1=2\\
	 &u_{n+1}=2u_n
	\quad (n \ge 1)}$.  Công thức số hạng tổng quát của dãy số là
	\choice 
		{$u_n=n^{n-1}$} 
		 {\True $u_n=2^n$}
		{$u_n=2^{n+1}$} 
		{$u_n=2$} 
	\loigiai{ 
		Ta có $\heva{&u_1=2\\ 
		&u_2=2u_1\\ 
		 &u_3=2u_2\\
		&\cdots  \\
		&u_n=2u_{n-1}.}$  \\
		Nhân hai vế ta được $u_1\cdot u_2\cdot u_3\ldots  u_n=2\cdot 2^{n-1} \cdot u_1\cdot u_2\ldots u_{n-1} \Leftrightarrow u_n=2^n$.
	 }
\end{ex}

% nguồn: Ok_Toan11_TracNghiem_Lan2
\begin{ex}%[1D2V1-2]%[ma:0068da51]
	Cho dãy số $(u_n)$  với $\heva{&u_1=-2\\
	 &u_{n+1}=-2-\dfrac{1}{u_n}
	\quad (n  \ge 1)}$. Công thức số hạng tổng quát của dãy số này là
	 \choice
		{$u_n=-\dfrac{n-1}{n}$} 
		 {$u_n=\dfrac{n+1}{n}$}
		{\True  $u_n=-\dfrac{n+1}{n}$}
		{$u_n=-\dfrac{n}{n+1}$} 
	\loigiai{Ta có $u_1=-\dfrac{3}{2}$; $u_2=-\dfrac{4}{3}$; $u_3=-\dfrac{5}{4}$; $\ldots$ Dễ dàng  dự đoán được $u_n=-\dfrac{n+1}{n}$
	 }
\end{ex}

% nguồn: Ok_Toan11_TracNghiem_Lan2
\begin{ex}%[1D2V1-2]%[ma:3a783a7a]
	Cho dãy số $(u_n)$ bởi công thức truy hồi sau $\heva{&u_1=0\\&u_{n+1}=u_n+n, \forall n \ge 1}$. Tính giá trị của  $u_{218}$.
	\choice 
		 {\True $23653$}
		 {$46872$}
		 {$23871$}
		 {$23436$}
	\loigiai{Ta có  $u_2=u_1+1=0+1=1$.\\
		 $u_3=u_2+2=1+2=3$.\\
		 $u_4=u_3+3=3+3=6$.\\
		$u_5=u_4+4=6+4=10$.\\ 
		$u_6=u_5+5=10+5=15$.\\ 
		 Dựa vào các số hạng trên, ta dự đoán công thức số hạng tổng quát của dãy $(u_n)$ là $u_n=\dfrac{1}{2}(n^2-n)$.\\
		Do  đó $u_{218}=\dfrac{1}{2}(218^2-218)=23653$.
	 }
\end{ex}

% nguồn: Ok_Toan11_TracNghiem_Lan2
\begin{ex}%[1D2V1-2]%[ma:9e6bbbaf]
	Cho dãy số $(u_n)$ với  $\heva{&u_1 = 1 \\
	&u_{n+1} = u_n  + (-1)^{2n+1}, \ n \in \mathbb{N}^*.}$\\
	Số  hạng tổng quát $u_n$ của dãy số này là số hạng nào dưới đây?
	 \choice
		{\True $u_n = 2 - n$} 
		 {$u_n$ không xác định}
		{$u_n  = 1 - n$}
		{$u_n = -n$  với mọi $n$}
	 \loigiai{
		Ta  có $(-1)^{2n+1} = -1$ cho mọi $n$. Do đó, công thức truy hồi trở thành $u_{n+1} = u_n - 1$. Khi đó
		\begin{align*}
			u_1  &= 1, \\
			u_2 -  u_1 &= -1, \\
			u_3 - u_2 &= -1,  \\
			u_4 - u_3 &=  -1, \\
			&\ldots \\ 
			 u_n - u_{n-1} &= -1.
		\end{align*}
		Cộng vế theo vế các đẳng thức trên, ta  được:
		 \[
		u_n = 1 +  (-1)(n-1) = 2 - n.
		\] 
		Vậy số hạng tổng quát của dãy số là  $u_n = 2 - n$.
	} 
\end{ex}

% nguồn: Ok_Toan11_TracNghiem_Lan2
\begin{ex}%[1D2V1-2]%[ma:35c6d379]
	Cho dãy số $(u_n)$ với $\heva{&u_1=1  \\
	&u_{n+1}=u_n + (-1)^{2n},  \ n \in \mathbb{N}^*.}$\\
	Số hạng tổng quát $u_n$ của  dãy số này là số hạng nào dưới đây?
	 \choice
		{\True $u_n = 1 + n$} 
		{$u_n = 1 - n$} 
		 {$u_n = 1 + (-1)^{2n}$}
		{$u_n  = n$}
	\loigiai{ 
		Ta có $(-1)^{2n} = 1$ cho mọi $n \in \mathbb{N}^*$. Do đó công thức truy hồi trở thành $u_{n+1} = u_n +  1$. Khi đó
		\begin{align*}
			 u_1 &= 1, \\
			u_2 -  u_1 &= 1, \\
			u_3  - u_2 &= 1, \\
			&\ldots \\ 
			u_n - u_{n-1} &= 1. 
		\end{align*}
		Cộng  vế theo vế các đẳng thức trên, ta được:
		\[ 
		u_n = 1  + (1 + 1 + \ldots + 1) = 1 + (n - 1) = 1 + n.
		 \]
		Vậy số hạng tổng quát của dãy số là $u_n =  1 + n$.
	} 
\end{ex}

% nguồn: Ok_Toan11_TracNghiem_Lan2
\begin{ex}%[1D2V1-2]%[ma:b361c2b5]
	Cho  dãy số $(u_n)$ với $\heva{&u_1=\dfrac{1}{2} \\
	&u_{n+1}=u_n-2 \quad (n \ge 1)}$. Công thức số hạng tổng quát của dãy số  này là
	 \choice
		 {$u_n=\dfrac{1}{2}+2(n-1)$}
		{\True $u_n=\dfrac{1}{2}-2(n-1)$} 
		 {$u_n=\dfrac{1}{2}-2n$}
		{$u_n=\dfrac{1}{2}+2n$} 
	\loigiai{ 
		 Ta có $\heva{&u_1=\dfrac{1}{2} \\
		 &u_2=u_1-2\\
		&u_3=u_2-2\\ 
		&\cdots \\ 
		&u_n=u_{n-1}-2.}$  \\
		Cộng hai vế ta được  $u_n=\dfrac{1}{2}-2-2\ldots-2=\dfrac{1}{2}-2(n-1)$.
	 }
\end{ex}

% nguồn: Ok_Toan11_TracNghiem_Lan2
\begin{ex}%[1D2V1-2]%[ma:4c5dbb3b]
	Cho dãy số $(u_n)$ thỏa  $\heva{&u_1=3\\
	&u_n=2u_{n-1}+n^2-n-3, \forall n \in  \mathbb{N}, n \geq 2}$. \\
	Tính  tổng $S_{20}=u_1+u_2+\cdots+u_{20}$.
	 \choice
		{$2\,022$} 
		 {$8\,385\,080$}
		 {$2\,021$}
		{\True  $8\,385\,087$}
	\loigiai{ 
		Ta  có
		\allowdisplaybreaks 
		\begin{eqnarray*}
			& & u_n=2u_{n-1}+n^2-n-3\\ 
			&\Leftrightarrow & u_n+n^2+3n+1=2u_{n-1}+2n^2-4n+2+6n-6+2\\ 
			 &\Leftrightarrow & u_n+n^2+3n+1=2u_{n-1}+2(n-1)^2+6(n-1)+2\\
			&\Leftrightarrow  & u_n+n^2+3n+1=2\left(u_{n-1}+(n-1)^2+3(n-1)+1\right).
		\end{eqnarray*}
		\allowdisplaybreaks 
		\begin{eqnarray*}
			\forall n \in \mathbb{N}, n \geq 1, \text{đặt  } v_n=u_n+n^2+3n+1\Rightarrow\heva{&v_1=u_1+1^2+3\cdot 1+1=8\\
			&v_{n-1}=u_{n-1}+(n-1)^2+3(n-1)+1.} 
		\end{eqnarray*}
		Ta có $(v_n)\colon\heva{&v_1=8\\ 
		&v_n=2v_{n-1}, \forall n \in \mathbb{N}, n \geq 2}$ là một cấp số nhân với $v_1=8$,  công bội là $q=2$.\\
		Khi đó $v_n=8\cdot 2^{n-1}=2^{n+2}$.\\ 
		Suy ra  $u_n=2^{n+2}-n^2-3n-1$. \\
		 Vậy
		 \allowdisplaybreaks
		\begin{eqnarray*}
			S_n&=&u_1+\cdots+u_n\\ 
			&=&2^2\left(2+2^2+\cdots+2^n\right)-1-2^2-\cdots-n^2-3(1+2+\cdots+n)-n\\ 
			 &=&2^3\left(2^n-1\right)-\dfrac{n(n+1)(2n+1)}{6}-3\cdot \dfrac{n(n+1)}{2}-n.
		\end{eqnarray*}
		 Vậy $S_{20}=8\,385\,087$.
	 }
\end{ex}

% nguồn: Ok_Toan11_TracNghiem_Lan2
\begin{ex}%[1D2C1-2]%[ma:5f56bf57]
	 Tìm $x$ biết $\left(x+1\cdot2\right)+\left(x+2\cdot5\right)+\left(x+3\cdot8\right)+\cdots+\left(x+10\cdot29\right)=1\,200$.
	 \choice
		{$x=7$} 
		 {$x=8$}
		{$x=9$} 
		{\True  $x=10$}
	\loigiai{ 
		Ta  có
		\[\begin{aligned}
		&\left( x+1\cdot2 \right)+\left(  x+2\cdot5 \right)+\left( x+3\cdot8 \right)+\cdots+\left( x+10\cdot29 \right)=1\,200 \\
		\Leftrightarrow&10x+\left( 1\cdot2+2\cdot5+3\cdot8+\cdots+10\cdot29 \right)=1\,200\\ 
		 \Leftrightarrow&10x+1\,100=1\,200\Leftrightarrow x=10.
\end{aligned}\]}
\end{ex}

% nguồn: Ok_Toan11_TracNghiem_Lan2
\begin{ex}%[1D2C1-2]%[ma:61afbd7f]
	Tính tổng  $S=1\cdot2+3\cdot4+5\cdot6+\cdots+11\cdot12$.
	 \choice
		{\True $322$} 
		 {$321$}
		 {$320$}
		 {$319$}
	\loigiai{ 
		 Đặt $a_k=\left(2k-1\right)2k$, $k\in\mathbb{N}^*$ ta có
		\[\begin{aligned}
		 S_n&=1\cdot2+3\cdot4+5\cdot6+\cdots+\left(2n-1\right)2n=\sum\limits_{k=1}^na_k\\
		&=\sum\limits_{k=1}^n\left(2k-1\right)2k=4\sum\limits_{k=1}^nk^2-2\sum\limits_{k=1}^nk\\ 
		&=\dfrac{4n\left(n+1\right)\left(2n+1\right)}{6}-n\left(n+1\right)=\dfrac{n\left( n+1 \right)\left( 4n-1  \right)}{3}.
	\end{aligned}\]
 Vậy $S=1\cdot2+3\cdot4+5\cdot6+\cdots+11\cdot12=\dfrac{6\left( 6+1 \right)\left( 4\cdot6-1 \right)}{3}=322$.}
\end{ex}

% nguồn: Ok_Toan11_TracNghiem_Lan2
\begin{ex}%[1D2C1-2]%[ma:f80cd04d]
	 Tổng $S=2\cdot3+4\cdot5+6\cdot7+\cdots+20\cdot21$ có giá trị bằng
	 \choice
		{$1\,550$} 
		{$1\,655$} 
		 {\True $1\,650$}
		 {$1\,450$}
	\loigiai{ 
		Đặt $a_k=2k\left(2k+1\right)$, $k\in\mathbb{N}^*$ ta  có
		\[\begin{aligned}
		 S_n&=3\cdot4+5\cdot6+\cdots+2n\left(2n+1\right)=\sum\limits_{k=1}^na_k\\
		 &=\sum\limits_{k=1}^n2k\left(2k+1\right)=4\sum\limits_{k=1}^nk^2+2\sum\limits_{k=1}^nk\\
		&=\dfrac{4n\left(n+1\right)\left(2n+1\right)}{6}+n\left(n+1\right)=\dfrac{n\left(  n+1 \right)\left( 4n+5\right)}{3}.
	\end{aligned}\]
Vậy $S=3\cdot4+5\cdot6+\cdots+11\cdot12=\dfrac{10\left(10+1 \right)\left(4\cdot10+5 \right)}{3}=1\,650$.} 
\end{ex}

% nguồn: Ok_Toan11_TracNghiem_Lan2
\begin{ex}%[1D2C1-2]%[ma:890c81d2]
	Tính tổng  $S_n=1\cdot5+3\cdot7+5\cdot9+\cdots+\left(2n-1\right)\left(2n+3\right)$ khi $n=15$.
	 \choice
		{$5\,450$} 
		 {$5\,400$}
		 {\True $5\,395$}
		 {$5\,650$}
	 \loigiai{
		Đặt $a_k=(2k-1)(2k+3)$,  $k\in\mathbb{N}^*$ ta có
		\[\begin{aligned}
		 S_n&=1\cdot5+3\cdot7+5\cdot9+\cdots+\left(2n-1\right)\left(2n+3\right)=\sum\limits_{k=1}^na_k\\
		 &=\sum\limits_{k=1}^n(2k-1)(2k+3)=4\sum\limits_{k=1}^nk^2+4\sum\limits_{k=1}^nk-3n\\
		&=\dfrac{2n\left(n+1\right)\left(2n+1\right)}3+2n\left(n+1\right)-3n=\dfrac{2n(n+1)(2n+4)}{3}-3n. 
	\end{aligned}\]
	Với $n=15$ ta có  $S_{15}=\dfrac{30\cdot16\cdot34}{3}-45=5\,395$.
} 
\end{ex}

% nguồn: Ok_Toan11_TracNghiem_Lan2
\begin{ex}%[1D2C1-2]%[ma:9fd31e10]
	Cho tổng $S_n=1\cdot2+3\cdot4+5\cdot6+\cdots+\left(2n-1\right)2n$.  Tính giá trị của $S_{50}$.
	 \choice
		{\True $169\,150$} 
		{$155\,000$} 
		 {$165\,050$}
		 {$165\,000$}
	\loigiai{ 
		Đặt $a_k=(2k-1)2k$, $k\in\mathbb{N}^*$  ta có
		\[\begin{aligned}
		S_n&=1\cdot2+3\cdot4+5\cdot6+\cdots+\left(2n-1\right)2n=\sum\limits_{k=1}^na_k\\ 
		&=\sum\limits_{k=1}^n(2k-1)2k=4\sum\limits_{k=1}^nk^2-2\sum\limits_{k=1}^nk\\ 
		 &=\dfrac{2n\left(n+1\right)\left(2n+1\right)}3-n\left(n+1\right)=\dfrac{n(n+1)(4n-1)}{3}.
	\end{aligned}\]
Vậy  $S_{50}=\dfrac{50\cdot51\cdot199}{3}=169\,150$.}
\end{ex}

% nguồn: Ok_Toan11_TracNghiem_Lan2
\begin{ex}%[1D2C1-2]%[ma:1c0c64fc]
	Giá trị của tổng $S_n=1\cdot4+3\cdot8+5\cdot12+\cdots+\left(2n-1\right)\cdot4n$  khi $ n=10$ là
	 \choice
		{$1\,650$} 
		 {\True $2\,860$}
		 {$2\,650$}
		 {$1\,950$}
	 \loigiai{
		 Đặt $a_k=(2k-1)4k$, $k\in\mathbb{N}^*$ ta có
		\[\begin{aligned}
		 S_n&=1\cdot4+3\cdot8+5\cdot12+\cdots+\left(2n-1\right)\cdot4n=\sum\limits_{k=1}^na_k\\
		&=\sum\limits_{k=1}^n(2k-1)4k=8\sum\limits_{k=1}^nk^2-4\sum\limits_{k=1}^nk\\ 
		 &=\dfrac{4n\left(n+1\right)\left(2n+1\right)}3-2n\left(n+1\right)=\dfrac{2n(n+1)(4n-1)}{3}.
	\end{aligned}\]
Với $n=10$ ta  có $S_{10}=\dfrac{20\cdot11\cdot39}{3}=2\,860$.}
\end{ex}

% nguồn: Ok_Toan11_TracNghiem_Lan2
\begin{ex}%[1D2C1-2]%[ma:b55a6eb5]
	Giá trị của tổng $S=1\cdot2+2\cdot5+3\cdot8+\cdots+20\cdot59$ là 
	 \choice
		 {$8\,450$}
		{$8\,300$} 
		 {$8\,850$}
		{\True  $8\,400$}
	 \loigiai{
		Đặt $a_k=k\left(3k-1\right)$,  $k\in\mathbb{N}^*$ ta có
		\[\begin{aligned}
		 S_n&=1\cdot2+2\cdot5+3\cdot8+\cdots+n(3n-1)=\sum\limits_{k=1}^na_k\\
		&=\sum\limits_{k=1}^nk\left(3k-1\right)=3\sum\limits_{k=1}^nk^2-\sum\limits_{k=1}^nk\\ 
		&=\dfrac{n\left(n+1\right)\left(2n+1\right)}2-\dfrac{n\left(n+1\right)}2=n^2(n+1). 
	\end{aligned}\]
Vậy $S=1\cdot2+2\cdot5+3\cdot8+\cdots+20\cdot59=20^2(20+1)=8\,400$.} 
\end{ex}

% nguồn: Ok_Toan11_TracNghiem_Lan2
\begin{ex}%[1D2C1-2]%[ma:0eda8e8f]
	Cho tổng  $S_n=1+2+2^2+\cdots+2^n$. Chọn mệnh đề đúng:
	 \choice
		{\True  $S_{10}=2\,047$}
		{$S_{10}=2\,048$} 
		 {$S_{10}=1\,024$}
		{$S_{10}=1\,023$} 
	\loigiai{ 
		Ta có  $S_n=\dfrac{2^{n+1}-1}{2-1}=2^{n+1}-1$.\\
		 Do đó $S_{10}=2^{11}-1=2\,047$.
	 }
\end{ex}


%%%%%%%%%%%%%%%% ĐÚNG/SAI %%%%%%%%%%%%%%%%

% nguồn: Ok_dungsai_lop11_Lan1
\begin{ex}%[1D2H1-2]%[ma:d2c22ccb]
 Cho dãy số $(u_n)$ biết $u_1=2$ và  $u_n=\sqrt{2+u_{n-1}^2}$ với mọi $n\ge 2$.
 \choiceTF 
  {Năm số hạng  đầu của dãy số $(u_n)$ là $2$; $\sqrt{2}$; $2\sqrt{2}$; $2\sqrt{3}$; $2\sqrt{4}$}
  {\True Dãy số $(u_n)$  là dãy tăng}
  {\True Tích $6$ số hạng đầu  của dãy số là $S=96\sqrt{35}$}
  {\True Số hạng tổng  quát $u_n=\sqrt{2(n+1)}$}
 \loigiai{ 
  \begin{itemchoice}
   \itemch Sai. Năm số hạng đầu của dãy số $(u_n)$ là $2$;  $\sqrt{6}$; $2\sqrt{2}$; $\sqrt{10}$; $2\sqrt{3}$. Vậy mệnh đề ``Năm số hạng đầu của dãy số $(u_n)$ là $2$; $\sqrt{2}$; $2\sqrt{2}$; $2\sqrt{3}$; $2\sqrt{4}$'' là mệnh đề sai.
   \itemch Đúng. Dễ thấy $ u_n >0, \forall n \in \mathbb{N}^{*} $.\\ Lại có, với mọi  $ n \in \mathbb{N}^{*}$ ta có $2+u_{n-1}^2>u_{n-1}^2 \Rightarrow \sqrt{2+u_{n-1}^2}>\sqrt{u_{n-1}^2}\Rightarrow u_n>u_{n-1} $.\\
   Vậy dãy  số $(u_n)$ là dãy tăng.
   \itemch Đúng. Vì sáu số hạng đầu của  dãy số $(u_n)$ là $2$; $\sqrt{6}$; $2\sqrt{2}$; $\sqrt{10}$; $2\sqrt{3}$; $\sqrt{14}$.\\
   Tích $6$ số hạng  đầu của dãy số là $S=96\sqrt{35}$.
   \itemch Đúng. Thật  vậy\\
   Ta  có $u_1=\sqrt{2(1+1)}=2$ (đúng).\\
   Giả sử với $n=k$, ta có $u_k=\sqrt{2(k+1)}$.\\ 
    Xét $u_{k+1}=\sqrt{2(k+1+1)}=\sqrt{2(k+1)+2}=\sqrt{2+u_k^2}$ (đúng).\\
    Vậy công thức số hạng tổng quát $u_n=\sqrt{2(n+1)}$ là mệnh đề đúng.
  \end{itemchoice}
 } 
\end{ex}

% nguồn: Ok_dungsai_lop11_Lan1
\begin{ex}%[1D2H1-2]%[ma:f71b72d2]
 Cho dãy số $\left(u_n\right)$ cho bởi công thức $u_n = \dfrac{3n-1}{n+1}$. Xét tính đúng sai của  các khẳng định sau
 \choiceTF 
  {Năm số hạng  đầu là $1;\dfrac{5}{3};2;\dfrac{12}{5};\dfrac{7}{3}$}
  {\True Số hạng $u_{n+1} =  \dfrac{3n+2}{n+2}$}
  {Dãy số $\left(u_n\right)$ là  dãy giảm}
   {\True Dãy số $\left(u_n\right)$ bị chặn}
  \loigiai{
  \begin{itemchoice}
   \itemch \textbf{Sai}. Ta có $u_1 = \dfrac{3 \cdot 1 - 1}{1+1} = 1$, tương tự ta cũng có $u_2 =  \dfrac{5}{3}$, $u_3 = 2$, $u_4 = \dfrac{11}{3} \neq \dfrac{12}{3}$, $u_5 = \dfrac{7}{3}$.
    \itemch \textbf{Đúng}. Từ $u_{n} = \dfrac{3n-1}{n+1}$ nên $u_{n+1} = \dfrac{3\left(n+1\right) -1}{\left(n+1\right)+1} = \dfrac{3n+2}{n+2}$.
   \itemch \textbf{Sai}. Ta kiểm  tra dấu của hiệu số $u_{n+1}-u_n$.\\
   \[u_{n+1}-u_n = \dfrac{3n+2}{n+2}  - \dfrac{3n-1}{n+1} = \dfrac{4}{\left(n+1\right)\left(n+2\right)} >0.\]
   Nên  $\left(u_n\right)$ là dãy tăng.
   \itemch \textbf{Đúng}. Theo lí thuyết, để chứng minh dãy số bị chặn ta chứng minh nó bị chặn trên và bị chặn dưới. 
   Với $n \in \mathbb{N}^\ast$ ta  có $u_n = \dfrac{3n-1}{n+1} >0$.\\
   Nên dãy $\left(u_n\right)$ bị chặn dưới bởi $0$.\\ 
   Mặt khác $u_n = \dfrac{3n-1}{n+1} = \dfrac{3n+3-4}{n+1} = 3 - \dfrac{4}{n+1} <3$, $\forall n \in \mathbb{N}^\ast$.\\ 
   Nên dãy $\left(u_n\right)$ bị chặn trên bởi $3$.\\ 
   Vậy dãy số  $\left(u_n\right)$ bị chặn.
  \end{itemchoice}
  }
\end{ex}

% nguồn: Ok_dungsai_lop11_Lan1
\begin{ex}%[1D2H1-2]%[ma:9e88f53a]
 Cho dãy số $u_n=20\sin \dfrac{n\pi}{2}+21\cos\dfrac{n\pi}{3}$. Xét tính đúng sai các mệnh  đề sau
  \choiceTF
  {Dãy số được  cho bởi phương pháp truy hồi}
  {\True Ba số hạng  đầu của dãy số là $\dfrac{61}{2}$; $-\dfrac{21}{2}$; $-41$}
  {Tổng năm số hạng đầu  của dãy số bằng $0$}
  {\True  $u_{n+12}=u_n$}
  \loigiai{
  \begin{itemchoice}
   \itemch Dãy  được cho bởi công thức của số hạng tổng quát $u_n$.\\
   Vậy ``Dãy số được cho bởi phương  pháp truy hồi'' là mệnh đề sai.
    \itemch Ba số hạng đầu của dãy số là $\dfrac{61}{2}$; $-\dfrac{21}{2}$; $-41$.
   Vậy mệnh đề  ``Ba số hạng đầu của dãy số là $\dfrac{61}{2}$; $-\dfrac{21}{2}$; $-41$'' là mệnh đề đúng.
   \itemch Năm số hạng đầu của dãy số là $\dfrac{61}{2}$;  $-\dfrac{21}{2}$; $-41$; $-\dfrac{21}{2}$; $\dfrac{61}{2}$.\\
   Suy ra tổng năm  số hạng đầu của dãy số bằng $-1$.\\
   Vậy ''Tổng năm số hạng đầu của dãy số  bằng $0$'' là mệnh đề sai.
   \itemch Ta có 
   \begin{eqnarray*}
    u_{n+12}&=&20\sin \dfrac{(n+12)\pi}{2}+21\cos \dfrac{(n+12)\pi}{3}=20\sin \left(\dfrac{n\pi}{2}+\dfrac{12\pi}{2}\right)+21\cos  \left(\dfrac{n\pi}{3}+\dfrac{12\pi}{3}\right)\\
    &=&20\sin \left(\dfrac{n\pi}{2}+6\pi\right)+21\cos  \left(\dfrac{n\pi}{3}+4\pi\right)\\
     &=&20\sin \dfrac{n\pi}{2}+21\cos \dfrac{n\pi}{3}\\
     &=&u_n.
   \end{eqnarray*}
  \end{itemchoice}
  }
\end{ex}

% nguồn: Ok_dungsai_lop11_Lan1
\begin{ex}%[1D2H1-2]%[ma:75d9933b]
 Cho dãy số $(u_n)$ có  số hạng tổng quát $u_n=\dfrac{2 n+1}{n+2}$. Khi đó
 \choiceTF 
  {\True Số hạng  đầu tiên của dãy số là $1$}
  {\True Số hạng $u_2=\dfrac{5}{4}$; $u_3=\dfrac{7}{5}$} 
  {\True Dãy số $(u_n)$ tăng} 
  {Số $\dfrac{167}{84}$ là số  hạng thứ $252$ của dãy số $(u_n)$}
 \loigiai{ 
 \begin{itemchoice}
  \itemch \textbf{Đúng}. Vì $u_1=\dfrac{2\cdot 1+1}{1+2}=1$. 
   \itemch \textbf{Đúng}. Vì $u_2=\dfrac{2\cdot 2+1}{2+2}=\dfrac{5}{4}$; $u_3=\dfrac{2\cdot 3+1}{3+2}=\dfrac{7}{5}$.
  \itemch \textbf{Đúng}.  Vì $u_{n+1}-u_n=\dfrac{2(n+1)+1}{n+3}-\dfrac{2n+1}{n+2}=\dfrac{3}{(n+3)(n+2)}>0$.
  \itemch \textbf{Sai}. Xét $\dfrac{2 n+1}{n+2}=\dfrac{167}{84}  \Leftrightarrow 84(2 n+1)=167(n+2) \Leftrightarrow n=250$.
   \\
  Vậy $\dfrac{167}{84}$ là số hạng thứ $250$ của dãy số  $(u_n)$.
 \end{itemchoice}
  }
\end{ex}

% nguồn: Ok_dungsai_lop11_Lan1
\begin{ex}%[1D2H1-2]%[ma:49f2356a]
  Cho dãy số $(u_n)$ biết $u_1=2$ và $u_n=\sqrt{2+u_{n-1}^2}$ với mọi $n\ge 2$. Xét tính đúng sai các mệnh đề sau
  \choiceTF
  {Năm số hạng đầu của dãy số $(u_n)$ là $2$;  $\sqrt{2}$; $2\sqrt{2}$; $2\sqrt{3}$; $2\sqrt{4}$}
  {\True Dãy  số $(u_n)$ là dãy vô hạn}
  {\True Tích $6$ số hạng đầu  của dãy số là $S=96\sqrt{35}$}
  {\True Số hạng tổng quát $u_n=\sqrt{2(n+1)}$} 
  \loigiai{
  \begin{itemchoice}
   \itemch Năm số hạng đầu  của dãy số $(u_n)$ là $2$; $\sqrt{6}$; $2\sqrt{2}$; $\sqrt{10}$; $2\sqrt{3}$. Vậy mệnh đề ``Năm số hạng đầu của dãy số $(u_n)$ là $2$; $\sqrt{2}$; $2\sqrt{2}$; $2\sqrt{3}$; $2\sqrt{4}$'' là mệnh đề sai.
   \itemch Dãy số $(u_n)$ là  dãy vô hạn.
   \itemch Sáu số hạng đầu của dãy số $(u_n)$ là $2$; $\sqrt{6}$; $2\sqrt{2}$; $\sqrt{10}$; $2\sqrt{3}$;  $\sqrt{14}$.\\
   Tích $6$ số hạng đầu của dãy số là  $S=96\sqrt{35}$.\\
   Vậy  mệnh đề ``Tích $6$ số hạng đầu của dãy số là $S=96\sqrt{35}$'' là mệnh đề đúng.
   \itemch Ta có  $u_1=\sqrt{2(1+1)}=2$ (đúng).\\
   Giả sử  với $n=k$, ta có $u_k=\sqrt{2(k+1)}$.\\
   Xét $u_{k+1}=\sqrt{2(k+1+1)}=\sqrt{2(k+1)+2}=\sqrt{2+u_k^2}$  (đúng).\\
   Vậy công thức số hạng tổng quát $u_n=\sqrt{2(n+1)}$ là mệnh đề  đúng.
  \end{itemchoice}
  }
\end{ex}

% nguồn: Ok_dungsai_lop11_Lan2
\begin{ex}%[1D2H1-2]%[ma:2e91932b]
	
	Xác định  tính đúng, sai của các khẳng định sau:
	 \choiceTF
		{\True Dãy số $\left(u_n\right)$ với $\dfrac{-2}{3}$,  $\dfrac{-1}{3}$; $0$, $\dfrac{1}{3}$, $\dfrac{2}{3}$, $1$, $\dfrac{4}{3}$ là cấp số cộng với $u_1=\dfrac{-2}{3}$, $d=\dfrac{1}{3}$}
		 {\True Dãy số $\left(u_n\right)$ với $u_n=7-3n$ là cấp số cộng với $u_1=4;d=-3$}
		{Dãy số $\left(u_n\right)$ với  $u_n=n^2+n+1$ là cấp số cộng với $u_1=3$, $d=1$}
		{\True Dãy số $\left(u_n\right)$ với $u_n=(-1)^n+3n$  không là cấp số cộng}
	 \loigiai{
		\begin{itemchoice}
			 \itemch \textbf{Đúng.}\\
			Dãy số $\left(u_n\right)$ với $\dfrac{-2}{3}$, $\dfrac{-1}{3}$; $0$, $\dfrac{1}{3}$, $\dfrac{2}{3}$,  $1$, $\dfrac{4}{3}$.\\
			Ta thấy  $u_2-u_1=u_3-u_2=u_4-u_3=\ldots=\dfrac{1}{3}$.\\
			Vậy $\left(u_n\right)$ là cấp  số cộng với $u_1=\dfrac{-2}{3}$, $d=\dfrac{1}{3}$.
			 \itemch \textbf{Đúng.}\\
			Dãy số $\left(u_n\right)$ với  $u_n=7-3n$.\\
			Ta có  $u_{n+1}-u_n=[7-3(n+1)]-(7-3n)=-3$.\\
			Vậy $\left(u_n\right)$ là cấp số  cộng với $u_1=7-3.1=4;d=-3$.
			\itemch  \textbf{Sai.}\\
			Dãy số  $\left(u_n\right)$ với $u_n=n^2+n+1$.\\
			 Ta có $u_{n+1}-u_n=(n+1)^2+(n+1)+1-\left(n^2+n+1\right)=2n+2$ phụ thuộc vào $n$.\\
			Vậy $\left(u_n\right)$ không là cấp  số cộng.
			\itemch \textbf{Đúng.}\\ 
			 Dãy số $\left(u_n\right)$ với $u_n=(-1)^n+3n$ .\\
			Ta có $u_{n+1}-u_n=(-1)^{n+1}+3(n+1)-\left[(-1)^n+3n\right]=-(-1)^n+3-(-1)^n=3-2(-1)^n$ phụ  thuộc vào $n$.\\
			Vậy $\left(u_n\right)$ không là cấp số cộng. 
		\end{itemchoice}
	} 
\end{ex}

% nguồn: Ok_dungsai_lop11_Lan2
\begin{ex}%[1D2H1-2]%[ma:4fbb49ad]
	 Cho dãy số $(u_n)$ xác định bằng hệ trức truy hồi: $u_1=1$, $u_{n+1}=u_n+\left( n+1 \right)$.
	\choiceTF 
		 {\True $u_3=6$}
		{\True $(u_n)$  là dãy số tăng}
		{Biết  rằng $1+2+\cdots+n=\dfrac{n\left( n+1 \right)}{2}$. Khi đó công thức của số hạng tổng quát là $u_{n+1}=\dfrac{\left( n+2 \right)\left( n+3 \right)}{2}$}
		{\True Tổng của hai số hạng liên tiếp  là một số chính phương}
	 \loigiai{
		\begin{itemchoice}
			\itemch \textbf{Đúng.}\\ 
			Ta  có $u_1=1$; $u_2=u_1+2=1+2=3$; $u_3=u_2+3=3+3=6$.
			\itemch \textbf{Đúng.}\\ 
			Ta có $u_{n+1}-u_n=n+1>0$,  $\forall n\in \mathbb{N}^*$ nên $(u_n)$ là dãy số tăng.
			\itemch \textbf{Sai.}\\ 
			Từ câu a), ta  nhận thấy
			$$u_{n+1}=1+2+\cdots+n+\left( n+1  \right)=\dfrac{n\left( n+1 \right)}{2}+\left( n+1 \right)=\dfrac{\left( n+1 \right)\left( n+2 \right)}{2}.$$
			\itemch \textbf{Đúng.}\\ 
			Theo  câu c) ta có
			$$u_{n+1}+u_n=\dfrac{\left( n+1 \right)\left( n+2 \right)}{2}+\dfrac{n\left( n+1 \right)}{2}=\dfrac{\left( n+1 \right)\left( 2n+2 \right)}{2}={{\left( n+1  \right)}^{2}}.$$
		\end{itemchoice}
	} 
\end{ex}

% nguồn: Ok_dungsai_lop11_Lan1
\begin{ex}%[1D2V1-2]%[ma:1ec8c70f]
 Cho  dãy số $(u_{n})$, biết $\heva{&u_1=-1\\&u_{n+1}=u_{n}+3}$ với $n\geq 1$.\\
 Xét tính đúng sai của các  khẳng định sau:
  \choiceTF
  {\True Bốn số hạng đầu tiên  của dãy số lần lượt là $-1$; $2$; $5$; $8$}
  {Số hạng thứ năm của dãy  là $13$}
  {Công thức số hạng tổng quát của  dãy số là: $u_{n} = 2n-3$}
   {\True $101$ là số hạng thứ $35$ của dãy số đã cho}
  \loigiai{
  \begin{itemchoice}
    \itemch Đúng.\\
   Ta có $u_{1}=-1$; $u_{2}=u_{1}+3=2$;  $u_{3}=u_{2}+3=5$; $u_{4}=u_{3}+3=8$.
   \itemch Sai.\\ $u_{5}=u_{4}+3=11$. 
   \itemch Sai.\\  Từ giả thiết, ta có $\heva{&u_1=-1\\&u_{2}=u_{1}+3\\&u_{3}=u_{2}+3\\& \ldots\\& u_{n}=u_{n-1}+3}$.\\
   Cộng theo vế toàn bộ các đẳng thức trên và triệt tiêu  các số hạng giống nhau ở hai vế, ta có $u_{n}=-1+3(n-1)=3n-4$.\\
   Vậy công thức số hạng tổng quát của dãy số là  $u_{n}=3n-4$.
   \itemch Đúng.\\ Xét $101=3n-4 \Rightarrow n=35$.\\ 
   Vậy $101$ là số hạng thứ $35$ của dãy  số đã cho.
 \end{itemchoice}}
\end{ex}

% nguồn: Ok_dungsai_lop11_Lan1
\begin{ex}%[1D2V1-2]%[ma:952e985b]
 Cho dãy số $(u_n)$,  biết $\heva{&u_1=-1\\&u_{n+1}=u_{n}+3}$ với $n\geq 1$.\\
  Xét tính đúng sai của các khẳng định sau:
  \choiceTF
  {\True Bốn số hạng đầu tiên của dãy số lần lượt là $-1$; $2$;  $5$; $8$}
  {Số hạng thứ  năm của dãy là $13$}
   {Công thức số hạng tổng quát của dãy số là $u_{n} = 3n-1$}
  {\True $101$ là số hạng thứ $35$ của dãy số đã  cho}
  \loigiai{
  \begin{itemchoice}
   \itemch \textbf{Đúng}.\\ 
   Ta có  $u_{1}=-1$; $u_{2}=u_{1}+3=2$; $u_{3}=u_{2}+3=5$; $u_{4}=u_{3}+3=8$.
   \itemch \textbf{Sai}.\\ $u_{5}=u_{4}+3=11$. 
    \itemch \textbf{Sai}.\\
   \textbf{Cách 1.} Ta có $u_{n+1}=u_{n}+3\Leftrightarrow u_{n+1}-u_{n}=3$ nên $(u_n)$ là cấp số cộng có số hạng đầu  $u_1=-1$ và công sai $d=3$.\\
   Công thức số hạng tổng quát của $(u_n)$  là $$u_{n}=-1+(n-1)\cdot3=3n-4.$$
    \textbf{Cách 2.}
   Từ  giả thiết, ta có $\heva{&u_1=-1\\&u_{2}=u_{1}+3\\&u_{3}=u_{2}+3\\& \ldots\\& u_{n}=u_{n-1}+3.}$\\
   Cộng theo vế toàn bộ các đẳng thức trên và triệt tiêu các số hạng giống nhau ở hai vế,  ta có $$u_{n}=-1+3(n-1)=3n-4.$$\\
   Vậy công thức số hạng tổng quát của dãy số là $u_{n}=3n-4$. 
   \itemch \textbf{Đúng}.\\ Xét $101=3n-4  \Rightarrow n=35$.\\
   Vậy $101$ là số hạng thứ  $35$ của dãy số đã cho.
 \end{itemchoice}}
\end{ex}

% nguồn: Ok_dungsai_lop11_Lan1
\begin{ex}%[1D2V1-2]%[ma:31c32b9d]
  Cho dãy số $(u_{n})$ biết $\heva{&u_1=2\\&u_{n+1}=u_{n}+5}$ với $n\geq 1$.\\
 Xét tính đúng sai của  các khẳng định sau:
 \choiceTF 
  {\True Năm số hạng đầu của dãy  số là $u_{1}=2$; $u_{2}=7$; $u_{3}=12$; $u_{4}=17$ và $u_{5}=22$}
  {\True Số hạng tổng quát  của dãy $(u_{n})$ là $u_{n}=5n-3$}
  {\True Số hạng  $u_{50}$ bằng $247$}
  {$512$ là số hạng thứ  $102$ của dãy}
  \loigiai{
  \begin{itemchoice}
   \itemch Đúng.\\ Ta có $u_{1}=2$; $u_{2}=u_{1}+5=7$; $u_{3}=u_{2}+5=12$;  $u_{4}=u_{3}+5=17$ và $u_{5}=u_{4}+5=22$.
   \itemch Đúng.\\ Ta có $\heva{&u_1=2\\&u_{2}=u_{1}+5\\&u_{3}=u_{2}+5\\&  \ldots\\& u_{n}=u_{n-1}+5}$.\\
   Cộng  theo vế toàn bộ các đẳng thức trên và triệt tiêu các số hạng giống nhau ở hai vế, ta có $u_{n}=2+5(n-1)=5n-3$.\\
   Vậy công thức số hạng tổng quát của  dãy số là $u_{n}=5n-3$.
   \itemch Đúng.\\ $u_{50}=5\cdot 50-3=247$. 
   \itemch Sai.\\ Xét  $5n-3=512 \Rightarrow n=103$.\\
   Vậy $512$ là số  hạng thứ $103$ của dãy số $(u_{n})$.
 \end{itemchoice}}
\end{ex}

% nguồn: Ok_dungsai_lop11_Lan1
\begin{ex}%[1D2V1-2]%[ma:0214e51e]
 Cho dãy số $(u_n)$ được  xác định như sau: $\heva{&u_1=2\\&u_{n+1}=u_n+5.}$ Khi đó
 \choiceTF 
   {\True Năm số hạng đầu của dãy số là $u_1=2$; $u_2=7$; $u_3=12$; $u_4=17$; $u_5=22$}
  {\True Số hạng tổng quát của dãy $(u_n)$  là $u_n=5n-3$}
   {\True Số hạng $u_{50}$ bằng $247$}
  {$512$  là số hạng thứ $102$ của dãy $(u_n)$}
 \loigiai{ 
  \begin{itemchoice}
    \itemch {\bf Đúng.}\\
   Năm số hạng đầu của dãy số là: $u_1=2$; $u_2=7$;  $u_3=12$; $u_4=17$; $u_5=22$.
   \itemch  {\bf Đúng.}\\
   Ta  có: $\heva{&u_1=2\\&u_2=u_1+5\\&u_3=u_2+5\\&u_4=u_3+5\\&\cdots \cdots \cdots \cdots \cdots\\&u_n=u_{n-1}+5.}$\\
   Cộng theo vế toàn  bộ đẳng thức trên rồi triệt tiêu các số hạng giống nhau ở hai vế, ta được $u_n=2+(n-1)5=5n-3$.
   \itemch {\bf Đúng.}\\ 
   Số hạng thứ 50 của dãy số là $u_{50}=5\cdot 50-3=247$. 
   \itemch {\bf Sai.}\\ 
   Xét $5n-3=512\Rightarrow n=103$.\\ 
   Vậy số  $512$ là số hạng thứ 103 của dãy số $(u_n)$.
  \end{itemchoice}
  }
\end{ex}

% nguồn: Ok_dungsai_lop11_Lan1
\begin{ex}%[1D2V1-2]%[ma:336a2412]
 Cho dãy  số $(u_{n})$ biết $u_{n}=\dfrac{1}{1\cdot 2}+\dfrac{1}{2\cdot 3}+ \dfrac{1}{3\dot 4}+\cdots+\dfrac{1}{n(n+1)}$.\\
 Xét tính đúng sai của các  khẳng định sau:
  \choiceTF
  {\True Số  hạng $u_{1}=\dfrac{1}{2}$}
  {\True Số hạng  $u_{3}=\dfrac{3}{4}$}
   {$\dfrac{10}{11}$ là số hạng thứ $11$ của dãy số}
  {$u_{2023}+ u_{2024}  >2$}
 \loigiai{ 
  \begin{itemchoice}
   \itemch Đúng.\\ Ta  có\\
   $$\heva{&\dfrac{1}{1\cdot 2}=\dfrac{1}{1}-\dfrac{1}{2}\\&\dfrac{1}{2\cdot 3}=\dfrac{1}{2}-\dfrac{1}{3}\\&\dfrac{1}{3\cdot 4}=\dfrac{1}{3}-\dfrac{1}{4}\\& \ldots\\&  \dfrac{1}{n\cdot (n+1)}=\dfrac{1}{n}-\dfrac{1}{n+1}}.$$\\
   Suy ra 
   \begin{eqnarray*}
    & u_{n}& =\dfrac{1}{1\cdot 2}+\dfrac{1}{2\cdot 3}+ \dfrac{1}{3\dot 4}+\cdots+\dfrac{1}{n(n+1)}\\ 
    &  &=\dfrac{1}{1}-\dfrac{1}{2} + \dfrac{1}{2}-\dfrac{1}{3} + \dfrac{1}{3}-\dfrac{1}{4} + \cdots -\dfrac{1}{n} + \dfrac{1}{n}-\dfrac{1}{n+1}\\
    & &=1-\dfrac{1}{n+1}\\ 
    &&=\dfrac{n}{n+1}. 
   \end{eqnarray*}
   Vậy số hạng tổng quát của  dãy số là $u_{n}=\dfrac{n}{n+1}$.\\
    Vậy $u_{1}=\dfrac{1}{2}$.
   \itemch Đúng.\\ Ta có  $u_{n}=\dfrac{n}{n+1}$.\\
   Vậy  $u_{3}=\dfrac{3}{4}$.
   \itemch Sai.\\ $\dfrac{10}{11}$ là số hạng thứ $10$  của dãy số.
   \itemch Sai.\\  $u_{2023}+ u_{2024} < 2$.
 \end{itemchoice}}
\end{ex}

% nguồn: Ok_dungsai_lop11_Lan1
\begin{ex}%[1D2V1-2]%[ma:76596a8c]
 Cho dãy số $(u_{n})$  biết $\heva{&u_{1}=2023; u_{2}=2024\\&2u_{n+1}=u_{n}+u_{n+2}}$ với $n\geq 1$.\\
  Xét tính đúng sai của các khẳng định sau:
 \choiceTF 
  {\True Dãy $(v_{n}): v_{n}=u_{n}-u_{n-1}$  là dãy không đổi}
  {\True Biểu thị $u_{n}$ qua $u_{n-1}$ ta được $u_{n}=u_{n-1}+1$} 
  {\True $u_{3}=2025$} 
  {$u_{2024}=4044$} 
 \loigiai{ 
  \begin{itemchoice}
   \itemch Đúng.\\ Ta có  $2u_{n+1}=u_{n}+u_{n+2} \Rightarrow u_{n+2}-u_{n+1}=u_{n+1}-u_{n} \Rightarrow v_{n+2}=v_{n+1}$.\\
   Tương tự, ta  chứng minh được $v_{n+1}=\cdots=v_2=1$ hay dãy $(v_{n})$ là dãy không đổi.
   \itemch Đúng.\\  Ta có $u_{n}-u_{n-1}=1 \Rightarrow u_{n}=u_{n-1}+1$.\\
   Suy  ra
   \begin{eqnarray*}
    & u_{n}& =(u_{n}-u_{n-1})+(u_{n-1}-u_{n-2})+ \cdots  + (u_{2}-u_{1})+u_{1}\\
    & &=1+1+\cdots+1+u_{1}\\ 
     & &=n-1+2023\\
     &&=n+2022.
   \end{eqnarray*}
   Khi đó $u_{2024}=4046$. 
   \itemch  Đúng.
   \itemch  Sai.
 \end{itemchoice}}
\end{ex}

% nguồn: Ok_dungsai_lop11_Lan1
\begin{ex}%[1D2V1-2]%[ma:ef16ac67]
 Cho dãy số $(u_{n})$ biết $\heva{&u_{1}=2023; u_{2}=2024\\&2u_{n+1}=u_{n}+u_{n+2}}$ với $n\geq 1$. Khi  đó
 \choiceTF 
  {\True Dãy  $(v_{n})\colon v_{n}=u_{n}-u_{n-1}$ là dãy không đổi}
  {\True Biểu  thị $u_{n}$ qua $u_{n-1}$ ta được $u_{n}=u_{n-1}+1$}
  {\True  $u_{3}=2025$}
  {$\lim\limits_{n\to  +\infty}\dfrac{u_n}{n}=2022$}
  \loigiai{
  \begin{itemchoice}
   \itemch\textbf{ Đúng.}\\ Ta có $2u_{n+1}=u_{n}+u_{n+2} \Rightarrow u_{n+2}-u_{n+1}=u_{n+1}-u_{n} \Rightarrow v_{n+2}=v_{n+1}$.\\ 
   Tương  tự, ta chứng minh được $v_{n+1}=\cdots=v_2=1$ hay dãy $(v_{n})$ là dãy không đổi.
   \itemch \textbf{Đúng.}\\ Ta  có $u_{n}-u_{n-1}=1 \Rightarrow u_{n}=u_{n-1}+1$.
    \itemch \textbf{Đúng.}\\
    Ta có
   \begin{eqnarray*}
    & u_{n}&  =(u_{n}-u_{n-1})+(u_{n-1}-u_{n-2})+ \cdots + (u_{2}-u_{1})+u_{1}\\
    &  &=1+1+\cdots+1+u_{1}\\
    &  &=n-1+2023\\
     &&=n+2022.
   \end{eqnarray*}
   Khi đó $u_{3}=2025$. 
    \itemch \textbf{Sai.}\\
    $\lim\limits_{n\to +\infty}\dfrac{u_n}{n}=\lim\limits_{n\to +\infty}\dfrac{n+2022}{n}=\lim\limits_{n\to +\infty }\left( 1+\dfrac{2022}{n}\right)=1$.
 \end{itemchoice}}
\end{ex}


%%%%%%%%%%%%%%%% TRẢ LỜI NGẮN %%%%%%%%%%%%%%%%

% nguồn: Ok_TraLoiNgan_Lop11_Lan1
\begin{ex}%[1D2N1-2]%[ma:b20444a9]
	 Tìm số hạng thứ 4 của dãy số $(u_n)$ xác định bởi$ \heva{&u_1=1 \\ &u_n=2u_{n-1}+3\quad\forall n\ge 2. }$
	 \shortans{29}
	 \loigiai{
		Ta có $u_1=1$; $u_2=2u_1+3=5 $; $u_3=2u_2+3=13$;  $u_4=2u_3+3=29$.
	 }
\end{ex}

% nguồn: Ok_TraLoiNgan_Lop11_Lan1
\begin{ex}%[1D2V1-2]%[ma:76a988c1]
	Cho dãy số $\left(u_{n} \right):-3;-1;1;3;5;\ldots$.  Tính số hạng thứ $7$.
	 \shortans{9}
	 \loigiai{
		Ta có $u_{1} =-3$, $u_{2}  =-1$, $u_{3} =1$, $u_{4} =3$, $u_{5} =5$, $\ldots $\\
		Suy  ra $u_{2} =u_{1} +2$, $u_{3} =u_{2} +2$, $u_{4} =u_{3} +2$, $u_{5} =u_{4} +2$, $\ldots $.\\
		Vậy hệ thức  truy hồi xác định dãy số đã cho là $\heva{&u_{1}=-3\\&u_{n+1}=u_{n}+2, (n \geq 1)}$.
	} 
\end{ex}

% nguồn: Ok_TraLoiNgan_Lop11_Lan1
\begin{ex}%[1D2V1-2]%[ma:1af5f2e4]
	Cho dãy số $(u_{n})$ với  $\heva{&u_{1}=2\\&u_{n+1}=2u_{n}, \forall n \geq 1}$. Tính $u_{13}$.
	 \shortans{8192}
	\loigiai{ 
		$u_{n+1}=2u_{n}, \forall n \geq 1 \Rightarrow \dfrac{u_{n+1}}{u_{n}}=2, \forall  n \geq 1$.\\
		Từ đó ta có $\dfrac{u_{2}}{u_{1}}=2$; $\dfrac{u_{3}}{u_{2}}=2$; $\dfrac{u_{4}}{u_{3}}=2$, $\ldots$;  $\dfrac{u_{n}}{u_{n-1}}=2$.\\
		Nhân theo vế tất cả các đẳng  thức trên ta được
		$$\dfrac{u_{n}}{u_{1}}=2^{n-1} \Leftrightarrow u_{n}=u_{1}\cdot  2^{n-1} \Leftrightarrow u_{n}=2\cdot 2^{n-1} \Leftrightarrow u_{n}=2^{n}.$$
	Do đó $u_{13}=8192$.} 
\end{ex}

% nguồn: Ok_TraLoiNgan_Lop11_Lan1
\begin{ex}%[1D2V1-2]%[ma:1217bcf0]
	\immini{Cho hình vuông $A_1 B_1 C_1 D_1$  có cạnh bằng $4$. Với mọi số nguyên dương $n\ge 2$, gọi $A_n, B_n, C_n, D_n$ lần lượt là trung điểm của các cạnh $A_{n-1} B_{n-1}, B_{n-1} C_{n-1}, C_{n-1} D_{n-1}$, $D_{n-1} A_{n-1}$. Gọi $S_n$ là diện tích của tứ giác $A_n B_n C_n D_n$. Tính giá trị của biểu thức $2^{25}\cdot S_{12}$.}
	{\begin{tikzpicture}[>=stealth,line join=round,line cap=round,font=\footnotesize,scale=0.8]
		 \path
			(0,0) coordinate (D_1) 
			(5,0)  coordinate (C_1)
			($(D_1)!1!90:(C_1)$) coordinate  (A_1)
			 ($(A_1)+(C_1)-(D_1)$) coordinate (B_1)
			($(A_1)!1/2!(B_1)$) coordinate  (A_2)
			($(A_1)!1/2!(D_1)$) coordinate (D_2) 
			($(D_1)!1/2!(C_1)$) coordinate (C_2) 
			($(B_1)!1/2!(C_1)$) coordinate  (B_2)
			 ($(A_2)!1/2!(B_2)$) coordinate (A_3)
			($(A_2)!1/2!(D_2)$)  coordinate (D_3)
			($(D_2)!1/2!(C_2)$) coordinate  (C_3)
			($(B_2)!1/2!(C_2)$) coordinate (B_3) 
			($(A_3)!1/2!(B_3)$) coordinate  (A_4)
			($(A_3)!1/2!(D_3)$) coordinate  (D_4)
			($(D_3)!1/2!(C_3)$) coordinate  (C_4)
			($(B_3)!1/2!(C_3)$) coordinate  (B_4);
		 \draw
			 (A_1)--(B_1)--(C_1)--(D_1)--(A_1)
			(A_2)--(B_2)--(C_2)--(D_2)--(A_2) 
			 (A_3)--(B_3)--(C_3)--(D_3)--(A_3)
			 (A_4)--(B_4)--(C_4)--(D_4)--(A_4);
		\foreach \x/\y in {A_1/180,B_1/0,C_1/-10,D_1/190,A_2/90,B_2/0,C_2/-90,D_2/180,A_3/60,B_3/-60,C_3/230,D_3/160,A_4/0,B_4/-90,C_4/180,D_4/90} 
		\draw[fill=black] (\x) circle (1.1pt) + (\y:0.5cm) node  {$\x$};
	\end{tikzpicture}
	} 
	\shortans{1024} 
	\loigiai{ 
		Ta thấy mỗi  tứ giác $A_n B_n C_n D_n$ là một hình vuông có cạnh là $A_n B_n$.\\
		Ta có $A_1 B_1=4,A_2 B_2=\dfrac{1}{\sqrt{2}} A_1 B_1=2\sqrt{2}, A_3  B_3=\dfrac{1}{\sqrt{2}} A_2 B_2=\left(\dfrac{1}{\sqrt{2}} \right)^2 A_1 B_1=2,\ldots$\\
		Tổng quát $A_n B_n=\left(\dfrac{1}{\sqrt{2}} \right)^{n-1} A_1 B_1=\left(\dfrac{1}{\sqrt{2}}  \right)^{n-1} \cdot 4=\left(\dfrac{1}{\sqrt{2}} \right)^{n+3}$.\\
		Diện tích hình vuông $A_n B_n C_n D_n$  là $S_n=\left(A_n B_n \right)^2=\left(\dfrac{1}{2} \right)^{n+3}$, với mọi $n\in \mathbb{N}^*$.\\
		Áp dụng với $n=12$ có  $S_{12}=\left(\dfrac{1}{2} \right)^{15}$.\\
		Vậy $2^{12}\cdot  S_{12}=2^{25}\cdot \dfrac{1}{2^{15}}=1024$
	} 
\end{ex}

% nguồn: Ok_TraLoiNgan_Lop11_Lan1
\begin{ex}%[1D2V1-2]%[ma:aa8c0da9]
	Cho dãy số  $(u_{n})$ với $\heva{&u_{1}=3\\&u_{n+1}=u_{n}+2, \forall n \geq 1}$ có số hạng tổng quát $u_{n}=an+b$ với $a,b\in \mathbb{Z}$. Tính $a+b$.
	\shortans{1} 
	\loigiai{ 
		$u_{n+1}=u_{n}+2 \Rightarrow u_{n+1}-u_{n}=2, \forall n \geq 1$. Khi đó ta  có
		$u_{2}-u_{1}=2$\\ 
		$u_{3}-u_{2}=2$\\ 
		 $\ldots$\\
		 $u_{n}-u_{n-1}=2$\\
		Cộng theo vế các đẳng thức  trên ta được
	$$u_{n}-u_{1}=(n-1) \cdot 2  \Leftrightarrow u_{n}=2n-1.$$}
\end{ex}

% nguồn: Ok_TraLoiNgan_Lop11_Lan2
\begin{ex}%[1D2V1-2]%[ma:bb7df0c9]
	Cho  dãy số $(u_n)$ xác định bởi $\heva{&u_1=1\\&u_{n+1}=u_n+2n+1,n\geq1}$. Giá trị của $n$ để\break $-u_n+2\,017n+2\,018=0$.
	 \shortans{2,02}
	\loigiai{ 
		Ta  có
		\begin{eqnarray*}
			&u_1&=1  \\
			 &u_2&=u_1+2\cdot1+1 \\
			&u_3&=u_3+2\cdot2+1  \\
			 &u_4&=u_3+2\cdot3+1 \\
			& &\ldots \\ 
			&u_n&=u_{n-1}+2(n-1)+1 
		\end{eqnarray*}
		 Cộng vế với vế các đẳng thức trên ta được
		 \[
		u_n=2(1+2+3+\cdots+n-1)+n=2\dfrac{(n-1+1)(n-1)}{2}+n=n^2,\forall  n\in\mathbb{N}^*.
		\] 
		Do đó 
		\[ 
		-u_n+2\,017n+2\,018=0\Leftrightarrow  -n^2+2\,017n+2\,018=0\Leftrightarrow\hoac{&n=-1~(\text{loại})\\&n=2\,018~(\text{nhận}).}
		\] 
		Vậy $n=2\,018$. 
	 }
\end{ex}


%%%%%%%%%%%%%%%% TỰ LUẬN %%%%%%%%%%%%%%%%

% nguồn: Ok_TuanLuan_Toan11_Lan2
\begin{ex}%[1D2H1-2]%[ma:0aac7b9c]
	Gọi $u_n$ là số đường tròn ở hàng thứ $n$ trong hình dưới đây. Dự đoán công thức  số hạng tổng quát của dãy số $(u_n)$. \\ [.2cm]
	\centerline{ 
	\begin{tikzpicture}[scale=0.7,>=stealth,line join=round,line cap=round,font=\footnotesize]
			\foreach  \x in {1,2,...,4}{
				\draw (0,{1-1.5*\x}) node [left] {Hàng  thứ \x};
				\foreach \y in {1,...,\x}{ 
				\draw ({1+1.3*(\y-1)},{1-1.5*\x}) circle  (.5);
			 }
			 }
		\end{tikzpicture}
	} 
	 \loigiai{
		Số đường tròn ở hàng thứ nhất là  $u_1=1$. \\
		Số đường tròn ở hàng thứ hai là $u_2=2$. \\ 
		 Số đường tròn ở hàng thứ ba là $u_3=3$. \\
		Số  đường tròn ở hàng thứ tư là $u_4=4$. \\
		Số đường tròn ở hàng thứ $n$ là $u_n=n$. \\ 
	 }
\end{ex}

% nguồn: Ok_TuanLuan_Toan11_Lan2
\begin{ex}%[1D2H1-2]%[ma:ebf799c9]
	Gọi $u_n$ là tổng diện tích các hình vuông có ở hàng thứ $n$ trong  hình bên dưới (mỗi ô vuông nhỏ là $1$ đơn vị diện tích). \\ [.2cm]
	 \centerline{
	\begin{tikzpicture}[scale=0.4,>=stealth,line join=round,line cap=round,font=\footnotesize]
			\foreach \x in  {1,2,...,4}{
				\draw  (0,{-.5*(\x)^2-1.5*(\x)+2+\x/2}) node [left] {Hàng thứ \x};
				\foreach \y in {1,...,\x}{ 
				\draw ({1+(\y-1)*(\x+1)},{-.5*(\x)^2-1.5*(\x)+2})  grid ++(\x,\x);
			} 
			} 
		\end{tikzpicture}
	 }
	\begin{listEX}[1]
		\item Tính $u_1$,  $u_2$, $u_3$, $u_4$.
		\item Dự đoán công thức số  hạng tổng quát của dãy số $(u_n)$.
	\end{listEX}
	 \loigiai{
		\begin{listEX}[1]
			\item Ta  có
			\begin{eqnarray*}
				&&u_1=1  \\
				&&u_2=2 \cdot 2^2=2^3=8.  \\
				&&u_3=3 \cdot 3^2=3^3=27.  \\
				&&u_4=4 \cdot  4^2=4^3=64.
			\end{eqnarray*}
			\item Từ các kết quả trên ta thấy $u_n=n^3$, $n\in  \mathbb{N^*}$.
		\end{listEX}
	} 
\end{ex}

% nguồn: Ok_TuanLuan_Toan11_Lan2
\begin{ex}%[1D2V1-2]%[ma:908badf2]
	Tính  $S=1+3+5+\cdots+4001$?
	\loigiai{ 
		\begin{itemize}
			\item Giải theo phương pháp tự luận 
			\begin{itemize}
				\item Ta có  $2S=(1+4001)+(3+3999)+(5+3997)+\cdots+(4001+1)=4002\cdot 2001$.
				\item $S=\dfrac{4002\cdot 2001}{2}=4004001$. 
			\end{itemize}
			\item Giải theo  phương pháp trắc nghiệm
			\begin{itemize}
				\item Số số hạng$n=\dfrac{4001-1}{2}+1=2001$. 
				\item Tổng$S=\dfrac{(1+4001)\cdot 2001}{2}=4004001$. 
			\end{itemize}
			 \item Giải theo Casio
			\begin{itemize}
				\item Công thức số hạng tổng quát của dãy là $u_n=u_1+(n-1)d=1+(n-1)\cdot 2=2n-1$. 
				\item  Số số hạng của dãy là $2001$.
				\item Nhập máy tính cho  ta kết quả$4004001$.
			\end{itemize}
			\item Những sai  lầm thường gặp
			\begin{itemize}
				\item Tính sai số số hạng của  dãy.
				\item Tìm sai công thức số hạng tổng quát của dãy số khi làm với máy  tính Casio.
			\end{itemize}
		\end{itemize}
		\begin{luuy}
			Nhận thấy việc tìm số hạng tổng quát của  dãy đối với HS trung bình, yếu là tương đối khó khăn. Vì thế ta nên sử dụng công thức giải nhanh để tìm số số hạng và tổng của dãy một cách nhanh chóng. Ở bài tập này thì việc vận dụng công thức tính nhẩm sẽ nhanh hơn Casio nhé các em!
		\end{luuy}
	} 
\end{ex}

% nguồn: Ok_TuanLuan_Toan11_Lan2
\begin{ex}%[1D2V1-2]%[ma:cebd627c]
	Cho tổng $S(n)=2+4+6+\cdots+2n$. Khi đó $S_{30}$ bằng? 
	\loigiai{ 
		\begin{itemize}
			 \item Giải theo phương pháp tự luận
			\begin{itemize}
				\item Ta có $S=2+4+6+\cdots+60$. 
				\item  $2S=(2+60)+(4+58)+(6+56)+\cdots+(60+2)$.
				\item $S(n)=\dfrac{(2+60)\cdot 30}{2}=930$. 
			\end{itemize}
			\item Giải theo  phương pháp trắc nghiệm
			\begin{itemize}
				 \item Số hạng thứ 30 $u_{30}=2\cdot 30=60$.
				\item  Số số hạng $n=30$.
				\item Tổng  $S=\dfrac{(2+60)\cdot 30}{2}=930$.
			\end{itemize}
			\item Giải theo  Casio
			\begin{itemize}
				\item Công thức số hạng tổng quát của  dãy là $2n$.
				\item Số  số hạng của dãy là $30$.
				\item Nhập máy tính  cho ta kết quả$930$.
			\end{itemize}
			\item Những sai lầm thường  gặp
			\begin{itemize}
				\item Tìm sai số hạng  thứ $n$.
			\end{itemize}
		\end{itemize}
		\begin{luuy}
			Trong  bài tập này HS cần chú ý tới số hạng tổng quát trong dãy đã cho sẵn. Từ đó sử dụng để tìm số hạng thứ n hoặc sử dụng trong việc bấm máy tính Casio một cách nhanh chóng tìm được kết quả.
		\end{luuy}
	 }
\end{ex}

% nguồn: Ok_TuanLuan_Toan11_Lan2
\begin{ex}%[1D2C1-2]%[ma:b9273dfc]
	Cho tổng $S=\dfrac{3}{(1\cdot 2)^2}+\dfrac{5}{(2\cdot  3)^2}+\dfrac{7}{(3\cdot 4)^2}+\cdots+\dfrac{2n+1}{[n(n+1)]^2}$. Tính $S_{10}$.
	 \loigiai{
		\begin{itemize}
			\item Giải theo  phương pháp tự luận
			\begin{itemize}
				\item  \textbf{Cách 1}
				 Ta có $\dfrac{3}{(1\cdot 2)^2}=\dfrac{1}{1^2}-\dfrac{1}{2^2};\dfrac{5}{(2\cdot 3)^2}=\dfrac{1}{2^2}-\dfrac{1}{3^2};\cdots$.
				Suy ra $S_n=\dfrac{1}{1^2}-\dfrac{1}{2^2}+\dfrac{1}{2^2}-\dfrac{1}{3^2}+\dfrac{1}{3^2}-\dfrac{1}{4^2}+\cdots+\dfrac{1}{n^2}-\dfrac{1}{(n+1)^2}=1-\dfrac{1}{(n+1)^2}$. 
				 Vậy$S_{10}=\dfrac{10(10+2)}{(10+1)^2}=\dfrac{120}{121}$.
				\item \textbf{Cách 2} 
				Ta có $S_{10}=\dfrac{3}{(1\cdot 2)^2}+\dfrac{5}{(2\cdot 3)^2}+\dfrac{7}{(3\cdot 4)^2}+\cdots+\dfrac{21}{(10\cdot 11)^2}$. 
				 Suy ra $S_{10}=\dfrac{1}{1^2}-\dfrac{1}{2^2}+\dfrac{1}{2^2}-\dfrac{1}{3^2}+\dfrac{1}{3^2}-\dfrac{1}{4^2}+\cdots+\dfrac{1}{10^2}-\dfrac{1}{11^2}=1-\dfrac{1}{11^2}=\dfrac{120}{121}$.
			\end{itemize}
			\item Giải theo phương  pháp trắc nghiệm
			Nhận thấy$\dfrac{3}{(1\cdot  2)^2}=\dfrac{1}{1^2}-\dfrac{1}{2^2};\dfrac{5}{(2\cdot 3)^2}=\dfrac{1}{2^2}-\dfrac{1}{3^2};\cdots$.
			Nhận định$S_n=1-\dfrac{1}{(n+1)^2}=\dfrac{n(n+2)}{(n+1)^2}$. 
			 Suy ra $S_{10}=\dfrac{10(10+2)}{(10+1)^2}=\dfrac{120}{121}$.
			 \item Giải theo phương pháp Casio
			Công thức  số hạng tổng quát của dãy là $u_n=\dfrac{2n+1}{[n(n+1)]^2}$.
			Số số hạng của dãy  là $n=10$.
			Nhập máy tính  cho ta kết quả$\dfrac{120}{121}$.
		\end{itemize}
		\begin{luuy}
			Học sinh cần chuyển các số  hạng của dãy về đúng dạng và tách theo công thức
			\[\dfrac{2na+a^2}{n^2(n+a)^2}=\dfrac{1}{n^2}-\dfrac{1}{(n+a)^2}.\] 
		\end{luuy}
	} 
\end{ex}

% nguồn: Ok_TuanLuan_Toan11_Lan2
\begin{ex}%[1D2C1-2]%[ma:30d14463]
	Tính tổng $S=4{,}5^{100}\cdot  \left( \dfrac{1}{5}+\dfrac{1}{5^2}+\dfrac{1}{5^3}+\cdots+\dfrac{1}{5^{100}} \right) +1$?
	\loigiai{ 
		\begin{itemize}
			\item Giải theo phương  pháp tự luận
			Đặt$M=1+\dfrac{1}{5}+\dfrac{1}{5^2}+\cdots+\dfrac{1}{5^{100}}$. 
			Ta có  $5M=5+1+\dfrac{1}{5}+\cdots+\dfrac{1}{5^{99}}$.
			\begin{align*}
				\Rightarrow 5M-M&=\left( 5+1+\dfrac{1}{5}+\cdots+\dfrac{1}{5^{99}}\right) - \left( 1+\dfrac{1}{5}+\dfrac{1}{5^2}+\cdots+\dfrac{1}{5^{100}}\right) =5-\dfrac{1}{5^{100}}\\ 
				 4M&=5-\dfrac{1}{5^{100}}\Rightarrow M=\dfrac{5}{4}-\dfrac{1}{4\cdot 5^{100}}\\
				S&=4{,}5^{100}\cdot \left( \dfrac{5}{4}-\dfrac{1}{4\cdot 5^{100}}\right) +1=5^{100} 
			\end{align*}
			\item Giải theo phương pháp  trắc nghiệm
			Đặt$M=\dfrac{1}{5}+\dfrac{1}{5^2}+\dfrac{1}{5^3}+\cdots+\dfrac{1}{5^{100}}$. 
			Áp dụng công thức tính nhanh với $u_1=\dfrac{1}{5}$; $a=\dfrac{1}{5}$ ta  có $M=\dfrac{\left( \dfrac{1}{5}\right) ^{100}-1}{\dfrac{1}{5}-1} =\dfrac{5^{100}-1}{4\cdot 5^{100}}$.
			Suy ra $  S=4{,}5^{100}\cdot\dfrac{5^{100}-1}{4\cdot 5^{100}}+1=5^{100}$.
		\end{itemize}
	} 
\end{ex}

% nguồn: Ok_TuanLuan_Toan11_Lan2
\begin{ex}%[1D2C1-2]%[ma:82d69a52]
	Cho tổng $S_n=\dfrac{1}{1\cdot 2\cdot 3}+\dfrac{1}{2\cdot 3\cdot  4}+\dfrac{1}{3\cdot 4\cdot 5}+\cdots+\dfrac{1}{n(n+1)(n+2)}$. Khi đó công thức của $S_n$ là
	\loigiai{ 
		\begin{itemize}
			\item Giải theo phương pháp tự  luận
			Ta có $\dfrac{2\cdot  1}{1\cdot 2\cdot 3}=\dfrac{1}{1\cdot 2}-\dfrac{1}{2\cdot 3};\dfrac{2\cdot 1}{2\cdot 3\cdot 4}=\dfrac{1}{2\cdot 3}-\dfrac{1}{3\cdot 4}$.
			Suy  ra
			\begin{align*}
				2S_n&=\dfrac{1}{1\cdot 2}-\dfrac{1}{2\cdot  3}+\dfrac{1}{2\cdot 3}-\dfrac{1}{3\cdot 4}+\cdots+\dfrac{1}{n(n+1)}-\dfrac{1}{(n+1)(n+2)}\\
				 &=\dfrac{1}{1\cdot 2}-\dfrac{1}{(n+1)(n+2)}=\dfrac{n(n+3)}{2(n+1)(n+2)}.
			\end{align*}
			Vậy  $S_n=\dfrac{n(n+3)}{4(n+1)(n+2)}$.
			\item Giải theo  phương pháp trắc nghiệm
			Nhận thấy $\dfrac{2\cdot 1}{1\cdot 2\cdot 3}=\dfrac{1}{1\cdot 2}-\dfrac{1}{2\cdot 3};\dfrac{2\cdot 1}{2\cdot  3\cdot 4}=\dfrac{1}{2\cdot 3}-\dfrac{1}{3\cdot 4};\cdots$.
			Nhận định $S=\dfrac{1}{1\cdot 2}-\dfrac{1}{(n+1)(n+2)}=\dfrac{n(n+3)}{4(n+1)(n+2)}$. 
			 \item Những sai lầm thường gặp.
			\begin{itemize}
				 \item Tách sai các số hạng.
			\end{itemize}
		\end{itemize}
		\begin{luuy}
			 Học sinh cần chuyển các số hạng của dãy về đúng dạng và tách theo công thức\[\dfrac{2a}{n(n+a)(n+2a)}=\dfrac{1}{n(n+a)}-\dfrac{1}{(n+a)(n+2a)}.\]
		\end{luuy}
	 }
\end{ex}

% nguồn: Ok_TuanLuan_Toan11_Lan2
\begin{ex}%[1D2C1-2]%[ma:82837c0c]
	Xác định công thức tính số hạng tổng quát $u_n$  theo $n$ của dãy số sau
	\[\heva{&u_1=1 \\&u_{n+1}=u_n+n^3}\quad \forall n  \geq 1.\]
	\loigiai{ 
		Ta có $u_{n+1}=u_n+n^3  \Rightarrow u_{n+1}-u_n=n^3$.
		Từ đó suy ra 
		 \[
		\begin{aligned}
			&  u_1=1 \\
			& u_2-u_1=1^3 \\ 
			 & u_3-u_2=2^3 \\
			 & u_4-u_3=3^3\\
			 &\cdots\\
			&  u_{n-1}-u_{n-2}=(n-2)^3 \\
			&  u_n-u_{n-1}=(n-1)^3.
		\end{aligned}
		 \]
		Cộng từng vế $n$ đẳng thức  trên
		 \[
		\begin{aligned}
			 & u_1+u_2-u_1+u_3-u_2+\cdots+u_{n-1}-u_{n-2}+u_n-u_{n-1}=1+1^3+2^3+3^3+\cdots+(n-2)^3+(n-1)^3 \\
			 \Leftrightarrow &u_n=1+1^3+2^3+3^3+\cdots+(n-2)^3+(n-1)^3.
		\end{aligned}
		\] 
		 Bằng phương pháp quy nạp ta chứng minh được$1^3+2^3+3^3+\cdots+(n-1)^3=\frac{(n-1)^2 \cdot n^2}{4}$.\\
		Vậy $u_n=1+\frac{n^2(n-1)^2}{4}$. 
	} 
\end{ex}

% nguồn: Ok_TuanLuan_Toan11_Lan2
\begin{ex}%[1D2C1-2]%[ma:52df9bb7]
	Cho dãy số $(a_n)$ có $a_k=\dfrac{1}{k(k+1)}$. Đặt  $u_n=\sum\limits_{k=1}^n a_k$. Xác định công thức tính $(u_n)$ theo $n$.
	 \loigiai{
		 Ta có $a_k=\dfrac{1}{k(k+1)}=\dfrac{1}{k}-\dfrac{1}{k+1}$, do đó
		\[u_n=\sum_{k=1}^n a_k=\left(1-\dfrac{1}{2}\right)+\left(\dfrac{1}{2}-\dfrac{1}{3}\right)+\cdots+\left(\dfrac{1}{n-1}-\dfrac{1}{n}\right)+\left(\dfrac{1}{n}-\dfrac{1}{n+1}\right)=1-\dfrac{1}{n+1}.\] 
	 }
\end{ex}

% nguồn: Ok_TuanLuan_Toan11_Lan2
\begin{ex}%[1D2C1-2]%[ma:8de015ec]
	Xác  định công thức tính số hạng tổng quát $u_n$ theo $n$ của dãy số sau $\heva{&u_1=3 \\ &u_{n+1}=u_n+2}$.
	 \loigiai{
		Ta  có
		\[ 
		\begin{aligned}
			& u_2=u_1+2=3+2=5  \\
			& u_3=u_2+2=5+2=7 \\ 
			& u_4=u_3+2=7+2=9 .  \\
			&  u_5=u_4+2=9+2=11 .
		\end{aligned}
		\] 
		Từ các số hạng đầu trên,  ta dự đoán số hạng tổng quát $u_n$ có dạng
		 \[
		u_n=2 n+1, \forall  n \geq 1(*)
		\] 
		Ta dùng phương pháp chứng minh quy nạp để chứng minh công thức (*) đúng. \\ 
		Với $n=1; u_1=2 \cdot 1+1=3$. Vậy $(*)$ đúng  với $n=1$.\\
		Giả sử $(*)$ đúng với $n=k$. Có nghĩa ta có $u_k=2  k+1$.\\
		Ta cần chứng minh $(*)$ đúng  với $n=k+1$. Có nghĩa là ta phải chứng minh
		 \[
		u_{k+1}=2(k+1)+1=2 k+3. 
		 \]
		Thật vậy  từ hệ thức xác định dãy số và theo (2) ta có
		\[ 
		u_{k+1}=u_k+2=2  k+1+2=2 k+3.
		\] 
		Vậy $(*)$ đúng khi $n=k+1$. Kết luận $(*)$ đúng với mọi số nguyên dương  $n$.
	} 
\end{ex}

% nguồn: Ok_TuanLuan_Toan11_Lan2
\begin{ex}%[1D2C1-2]%[ma:095fd471]
	Tính tổng sau $S=\dfrac{2}{1\cdot  3}+\dfrac{2}{3\cdot 5}+\dfrac{2}{5\cdot 7}+\cdots+\dfrac{2}{97\cdot 99}$.
	 \loigiai{
		\begin{itemize}
			\item Giải theo phương  pháp tự luận
			Ta có $\dfrac{1}{1\cdot 3}=\dfrac{1}{1}-\dfrac{1}{3};\dfrac{1}{3\cdot  5}=\dfrac{1}{3}-\dfrac{1}{5};\cdots$.
			Do  đó$S=\dfrac{1}{1}-\dfrac{1}{3}+\dfrac{1}{3}-\dfrac{1}{5}+\cdots+\dfrac{1}{97}-\dfrac{1}{99}=1-\dfrac{1}{99}=\dfrac{98}{99}$.
			\item Giải  theo phương pháp trắc nghiệm
			Nhận  thấy$\dfrac{1}{1\cdot 3}=\dfrac{1}{1}-\dfrac{1}{3};\dfrac{1}{3\cdot 5}=\dfrac{1}{3}-\dfrac{1}{5};\cdots$.
			 Nhận định$S=1-\dfrac{1}{99}=\dfrac{98}{99}$.
			\item Giải  bằng Casio
			 Số hạng tổng quát của dãy là $u_n=\dfrac{2}{(2n-1)(2n+1)}$.
			Số số hạng của dãy là $n=49(2n-1=97\Leftrightarrow  n=49)$.
			Nhập máy tính cho  ta kết quả$S=\dfrac{98}{99}$.
			\item Những sai lầm thường  gặp
			\begin{itemize}
				\item Tách sai các số  hạng.
				\item Tìm sai số hạng tổng quát của  dãy số.
			\end{itemize}
		\end{itemize}
		\begin{luuy}
			Học sinh cần chuyển các số hạng của dãy về  đúng dạng và tách theo công thức
			 \[\dfrac{a}{n(n+a)}=\dfrac{1}{n}-\dfrac{1}{n+a}.\]
			Ở bài tập này việc làm bằng máy tính Casio là khó khăn và phức  tạp hơn vì chưa có sẵn số hạng tổng quát và số số hạng.
		\end{luuy}
	 }
\end{ex}

% nguồn: Ok_TuanLuan_Toan11_Lan2
\begin{ex}%[1D2C1-2]%[ma:4ef88c95]
	Tính  tổng$S=1+3+3^2+\cdots+3^{50}$?
	 \loigiai{
		\begin{itemize}
			\item Giải theo phương  pháp tự luận
			Ta có $3S=3+3^2+3^3+\cdots+3^{51}$. 
			Suy  ra $ 3S-S=(3+3^2+3^3+\cdots+3^{51})-(1+3+3^2\cdots+3^{50})$.
			 Do đó $2S=3^{51}-1\Rightarrow S=\dfrac{3^{51}-1}{2}$.
			\item Giải theo phương pháp  trắc nghiệm
			Áp dụng công thức tính nhanh với $u_1=1$; $a=3$ ta có $S=\dfrac{3^{51}-1}{2}$ 
			\item Giải theo Casio 
			Công thức số hạng tổng quát của dãy là  $u_n=3^{n-1}$.
			Số số hạng  của dãy là $n=51$.
			 Nhập máy tính cho ta kết quả $1{,}076846982\cdot 10^{24}$.
			Ta gán$1{,}076846982\cdot 10^{24}\to  A$.
			Lấy từng kết quả ở 4 đáp án  trừ cho $A$ khi nào bằng $0$ thì chọn đáp án đó.
			\item Những sai lầm thường gặpTìm sai số hạng tổng  quát của dãy số.
		\end{itemize}
		\begin{luuy}
			Khi làm với máy tính Caiso các em cần tìm chính xác số hạng tổng quát của dãy số việc  này quyết định máy có đưa ra được kết quả chính xác hay không. Ở bài tập này nếu các em thuộc được công thức tính nhanh thì ta có thể giải quyết bài toán hết sức nhanh chóng.
			Chú ý rằng bài toán này có thể hạn chế Casio bằng cách cho 2 đáp án ở “gần nhau” chẳng hạn phương án B và $\dfrac{3^{51}}{2}$ thì khi làm bằng Casio sẽ có 2 đáp án không phân biệt được  là B và C.
		\end{luuy}
	} 
\end{ex}

% nguồn: Ok_TuanLuan_Toan11_Lan1
\begin{bt}%[1D2Y1-2]%[ma:99b65151]
	Trong mặt phẳng $(P)$ cho hình bình hành $ABCD$. Qua $A$, $B$, $C$, $D$ lần  lượt vẽ bốn đường thẳng $a$, $b$, $c$, $d$ đôi một song song với nhau và không nằm trong mặt phẳng $(P)$. Một mặt phẳng cắt $a$, $b$, $c$, $d$ lần lượt tại bốn điểm $A'$, $B'$, $C'$, $D'$. Chứng minh rằng $A'B'C'D'$ là hình bình hành.
	 \loigiai{
	\immini 
		{Giao tuyến của mặt phẳng $(\alpha)$ với các mặt phẳng $(Bb,Cc)$, $(Aa,Dd)$, $(Aa,Bb)$, $(Cc,Dd)$  lần lượt là $B'C'$, $A'D'$, $A'B'$, $C'D'$.\\
		Do $(Bb,Cc) \parallel (Aa,Dd)$, $(Aa,Bb) \parallel (Cc,Dd)$ nên $B'C'  \parallel A'D'$ và $A'B' \parallel C'D'$.\\
		Suy  ra tứ giác $A'B'C'D'$ là hình bình hành.
		}{\begin{tikzpicture}[scale=1,font=\footnotesize,line join=round,line cap=round,>=stealth]
				 \coordinate (A) at (0,0);
				\coordinate  (B) at (4.5,0);
				\coordinate (C)  at (3.2,-0.8);
				 \coordinate (D) at ($(A)+(C)-(B)$);
				\def\vx{0.5} \def\vy{4} 
				\coordinate (v)  at (\vx,\vy);
				\coordinate (atop) at ($(A)+(v)$); 
				\coordinate (btop)  at ($(B)+(v)$);
				\coordinate  (ctop) at ($(C)+(v)$);
				\coordinate (dtop)  at ($(D)+(v)$);
				 \coordinate (B') at ($(B)!0.52!(btop)$);
				 \coordinate (C') at ($(C)!0.86!(ctop)$);
				 \coordinate (D') at ($(D)!0.78!(dtop)$);
				\coordinate (A')  at ($(A)!0.44!(atop)$);
				\draw (D)--(C)--(B); 
				\draw (B)--(btop); 
				\draw (C)--(ctop); 
				 \draw (D)--(dtop);
				\draw (A')--(atop); 
				\draw  (B')--(C')--(D');
				 \draw[dashed] (B)--(A)--(D);
				\draw[dashed] (A)--(A'); 
				\draw[dashed] (B')--(A')--(D'); 
				\foreach \p in {A,B,C,D,A',B',C',D'} 
				\fill (\p)  circle (1pt);
				\node[shift={(150:0.3)}]  at (A){$A$};
				\node[right] at (B){$B$}; 
				\node[below] at (C){$C$}; 
				\node[below]  at (D){$D$};
				 \node[below left] at (A'){$A'$};
				 \node[right] at (B'){$B'$};
				\node[above  right] at (C'){$C'$};
				 \node[left] at (D'){$D'$};
				\node[below left] at (atop){$a$}; 
				 \node[below left] at (btop){$b$};
				\node[below left]  at (ctop){$c$};
				\node[below left] at  (dtop){$d$};
			\end{tikzpicture}
		 }
	} 
\end{bt}

% nguồn: Ok_TuanLuan_Toan11_Lan1
\begin{bt}%[1D2Y1-2]%[ma:53a84cc3]
	Cho tứ diện $ABCD$. Lấy $G_1$, $G_2$, $G_3$ lần lượt là trọng tâm của  các tam giác $ABC$, $ACD$, $ADB$
	\begin{enumerate}
		\item Chứng minh  rằng $\left(G_1G_2G_3\right)\parallel(BCD)$.
		\item Xác định giao tuyến của mặt  phẳng $\left(G_1G_2G_3\right)$ với mặt phẳng $(ABC)$.
	\end{enumerate}
	\loigiai{ 
	\immini{\begin{enumerate}
		\item Gọi $M,~N,~L$ lần lượt là trung điểm của các cạnh $BC,~CD,~BD$.\\ 
		Ta có $\dfrac{AG_1}{AM}=\dfrac{AG_2}{AN}=\dfrac{AG_3}{AL}=\dfrac{2}{3}$ suy ra $\left\{\begin{aligned}
		&G_1G_2 \parallel  MN \\
		 &G_2G_3 \parallel NL.
	\end{aligned}\right. $\\
	Mà $\left\{\begin{aligned}
	&MN\subset (BCD) \\ 
	&NL\subset (BCD) 
	\end{aligned}\right. $ nên $\left\{\begin{aligned}
	&G_1G_2 \parallel  (BCD) \\
	&G_2G_3 \parallel (BCD) 
	\end{aligned}\right. $\\
	suy ra  $(G_1G_2G_3) \parallel (BCD)$.
	\item Từ $(G_1G_2G_3) \parallel (BCD)$ suy ra  $BC \parallel (G_1G_2G_3)$.\\
	Ta có $\left\{\begin{aligned}
	 &G_1\in (ABC)\cap (G_1G_2G_3) \\
	 &BC\subset (ABC) \\
	&BC  \parallel (G_1G_2G_3)
	\end{aligned}\right.\\ \Rightarrow (ABC) \cap (G_1G_2G_3)=\Delta \parallel BC$ với $\Delta$ qua $G_1$.
	\end{enumerate}}{\begin{tikzpicture}[scale=1,line join=round,line cap=round]
		\tkzDefPoints{0/0/D,1.1/-1.7/B,4.5/0/C} 
		 \coordinate (A) at ($(D)+(1.5,4)$);
		 \coordinate (M) at ($(C)!0.5!(B)$);
		\coordinate (N) at ($(C)!0.5!(D)$); 
		\coordinate (L)  at ($(D)!0.5!(B)$);
		\coordinate (G_1) at ($(M)!1/3!(A)$); 
		 \coordinate (G_2) at ($(N)!1/3!(A)$);
		\coordinate  (G_3) at ($(L)!1/3!(A)$);
		\coordinate (E) at ($(B)!1/3!(A)$); 
		\coordinate (F) at ($(C)!1/3!(A)$); 
		\coordinate (I) at  ($(D)!1/3!(A)$);
		 \tkzDrawLines[add=1cm and 1cm](E,F)
		\coordinate  (M') at ($(E)!1.5!(F)$);
		\tkzDrawPolygon (A,D,B,C) 
		 \tkzDrawSegments (A,B A,M A,L)
		\tkzDrawSegments[dashed](M,N  N,L D,C A,N G_1,G_2 G_2,G_3 G_3,G_1)
		\tkzDrawPoints[fill=black](D,B,C,A,M,N,L,G_1,G_2,G_3) 
		 \tkzLabelPoints[above](A)
		 \tkzLabelPoints[above left](G_2)
		 \tkzLabelPoints[below](B,M,N)
		 \tkzLabelPoints[left](D,L)
		 \tkzLabelPoints[right](C)
		\tkzLabelPoints[below right](G_1) 
		\tkzLabelPoints[below left](G_3) 
		 \draw (M') node[above]{$\Delta$};
	\end{tikzpicture}
	}} 
\end{bt}

% nguồn: Ok_TuanLuan_Toan11_Lan1
\begin{bt}%[1D2Y1-2]%[ma:d2e2c2b4]
	Bạn Chung cho rằng:  Nếu mặt phẳng $(P)$ chứa hai đường thẳng $a$, $b$ và $a$, $b$ cùng song song với mặt phẳng $(Q)$ thì $(P)$ luôn song song với $(Q)$. Phát biểu của bạn Chung có đúng không? Vì sao?
	 \loigiai{
		Phát biểu của bạn Chung không đúng. Vì thiếu  giả thiết $a$, $b$ cắt nhau nên chúng có thể song song với nhau. Khi đó mặt phẳng $(P)$ và mặt phẳng $(Q)$ có thể cắt nhau theo một giao tuyến là đường thẳng $\Delta$ song song với $a$ và $b$.
	 }
\end{bt}

% nguồn: Ok_TuanLuan_Toan11_Lan2
\begin{ex}%[1D2Y1-2]%[ma:f4b29b78]
	Trong mặt phẳng $(P)$ cho hình bình hành $ABCD$. Qua $A$, $B$, $C$, $D$ lần lượt vẽ bốn đường thẳng $a$, $b$, $c$, $d$ đôi một song song với nhau và không nằm trong mặt phẳng $(P)$. Một mặt phẳng cắt $a$, $b$, $c$, $d$ lần lượt tại bốn điểm $A'$,  $B'$, $C'$, $D'$. Chứng minh rằng $A'B'C'D'$ là hình bình hành.
	\loigiai{ 
	 \immini
		{Giao tuyến của mặt phẳng $(\alpha)$ với các mặt phẳng $(Bb,Cc)$, $(Aa,Dd)$, $(Aa,Bb)$, $(Cc,Dd)$ lần lượt là $B'C'$, $A'D'$, $A'B'$,  $C'D'$.\\
		Do $(Bb,Cc) \parallel (Aa,Dd)$, $(Aa,Bb) \parallel (Cc,Dd)$ nên $B'C' \parallel A'D'$ và $A'B'  \parallel C'D'$.\\
		Suy ra  tứ giác $A'B'C'D'$ là hình bình hành.
		}{\begin{tikzpicture}[scale=1,font=\footnotesize,line join=round,line cap=round,>=stealth]
				\coordinate  (A) at (0,0);
				\coordinate (B)  at (4.5,0);
				 \coordinate (C) at (3.2,-0.8);
				\coordinate (D) at  ($(A)+(C)-(B)$);
				\def\vx{0.5}  \def\vy{4}
				\coordinate (v) at  (\vx,\vy);
				\coordinate (atop)  at ($(A)+(v)$);
				 \coordinate (btop) at ($(B)+(v)$);
				\coordinate (ctop) at  ($(C)+(v)$);
				\coordinate  (dtop) at ($(D)+(v)$);
				\coordinate (B') at  ($(B)!0.52!(btop)$);
				 \coordinate (C') at ($(C)!0.86!(ctop)$);
				\coordinate (D')  at ($(D)!0.78!(dtop)$);
				\coordinate (A') at ($(A)!0.44!(atop)$); 
				\draw  (D)--(C)--(B);
				\draw (B)--(btop); 
				\draw (C)--(ctop); 
				\draw  (D)--(dtop);
				\draw  (A')--(atop);
				 \draw (B')--(C')--(D');
				\draw[dashed] (B)--(A)--(D); 
				 \draw[dashed] (A)--(A');
				\draw[dashed] (B')--(A')--(D'); 
				\foreach \p in  {A,B,C,D,A',B',C',D'}
				\fill (\p)  circle (1pt);
				\node[shift={(150:0.3)}]  at (A){$A$};
				 \node[right] at (B){$B$};
				\node[below]  at (C){$C$};
				 \node[below] at (D){$D$};
				\node[below left] at (A'){$A'$}; 
				\node[right] at (B'){$B'$}; 
				 \node[above right] at (C'){$C'$};
				\node[left]  at (D'){$D'$};
				\node[below left]  at (atop){$a$};
				\node[below  left] at (btop){$b$};
				 \node[below left] at (ctop){$c$};
				\node[below left]  at (dtop){$d$};
			\end{tikzpicture}
		 }
	} 
\end{ex}

% nguồn: Ok_TuanLuan_Toan11_Lan2
\begin{ex}%[1D2Y1-2]%[ma:1570d1b3]
	Cho tứ diện $ABCD$. Lấy $G_1$, $G_2$, $G_3$ lần lượt là trọng tâm của các tam giác $ABC$,  $ACD$, $ADB$
	\begin{enumerate}
		\item Chứng minh rằng  $\left(G_1G_2G_3\right)\parallel(BCD)$.
		\item Xác định  giao tuyến của mặt phẳng $\left(G_1G_2G_3\right)$ với mặt phẳng $(ABC)$.
	\end{enumerate}
	 \loigiai{
	\immini{\begin{enumerate}
		\item Gọi $M,~N,~L$ lần lượt là trung điểm của các cạnh  $BC,~CD,~BD$.\\
		Ta có $\dfrac{AG_1}{AM}=\dfrac{AG_2}{AN}=\dfrac{AG_3}{AL}=\dfrac{2}{3}$ suy ra $\left\{\begin{aligned}
		&G_1G_2 \parallel MN  \\
		&G_2G_3  \parallel NL.
	\end{aligned}\right. $\\
	Mà $\left\{\begin{aligned}
	 &MN\subset (BCD) \\
	 &NL\subset (BCD)
	\end{aligned}\right. $ nên $\left\{\begin{aligned}
	&G_1G_2 \parallel (BCD) \\ 
	&G_2G_3 \parallel (BCD) 
	\end{aligned}\right. $\\
	suy ra $(G_1G_2G_3) \parallel  (BCD)$.
	\item Từ $(G_1G_2G_3) \parallel (BCD)$ suy  ra $BC \parallel (G_1G_2G_3)$.\\
	Ta có $\left\{\begin{aligned}
	&G_1\in (ABC)\cap (G_1G_2G_3) \\ 
	&BC\subset (ABC)  \\
	&BC \parallel  (G_1G_2G_3)
	\end{aligned}\right.\\ \Rightarrow (ABC) \cap (G_1G_2G_3)=\Delta \parallel BC$ với $\Delta$ qua $G_1$.
	\end{enumerate}}{\begin{tikzpicture}[scale=1,line join=round,line cap=round]
		\tkzDefPoints{0/0/D,1.1/-1.7/B,4.5/0/C} 
		\coordinate (A) at  ($(D)+(1.5,4)$);
		\coordinate (M) at ($(C)!0.5!(B)$); 
		\coordinate  (N) at ($(C)!0.5!(D)$);
		\coordinate  (L) at ($(D)!0.5!(B)$);
		 \coordinate (G_1) at ($(M)!1/3!(A)$);
		\coordinate (G_2) at ($(N)!1/3!(A)$); 
		\coordinate (G_3) at  ($(L)!1/3!(A)$);
		\coordinate (E) at  ($(B)!1/3!(A)$);
		\coordinate (F) at  ($(C)!1/3!(A)$);
		 \coordinate (I) at ($(D)!1/3!(A)$);
		\tkzDrawLines[add=1cm and  1cm](E,F)
		\coordinate  (M') at ($(E)!1.5!(F)$);
		\tkzDrawPolygon (A,D,B,C) 
		\tkzDrawSegments (A,B A,M  A,L)
		\tkzDrawSegments[dashed](M,N  N,L D,C A,N G_1,G_2 G_2,G_3 G_3,G_1)
		 \tkzDrawPoints[fill=black](D,B,C,A,M,N,L,G_1,G_2,G_3)
		\tkzLabelPoints[above](A) 
		\tkzLabelPoints[above left](G_2) 
		 \tkzLabelPoints[below](B,M,N)
		 \tkzLabelPoints[left](D,L)
		\tkzLabelPoints[right](C) 
		\tkzLabelPoints[below right](G_1) 
		 \tkzLabelPoints[below left](G_3)
		\draw  (M') node[above]{$\Delta$};
	\end{tikzpicture}
	 }}
\end{ex}

% nguồn: Ok_TuanLuan_Toan11_Lan2
\begin{ex}%[1D2Y1-2]%[ma:c50aecaa]
	Bạn Chung cho rằng: Nếu mặt phẳng $(P)$  chứa hai đường thẳng $a$, $b$ và $a$, $b$ cùng song song với mặt phẳng $(Q)$ thì $(P)$ luôn song song với $(Q)$. Phát biểu của bạn Chung có đúng không? Vì sao?
	\loigiai{ 
		Phát biểu của bạn Chung không đúng. Vì thiếu giả thiết $a$, $b$ cắt nhau nên chúng có  thể song song với nhau. Khi đó mặt phẳng $(P)$ và mặt phẳng $(Q)$ có thể cắt nhau theo một giao tuyến là đường thẳng $\Delta$ song song với $a$ và $b$.
	 }
\end{ex}
